\subsection{取值于李代数的微分形式的原函数}

引理 5. 设 X 是一个 \(\mathbf{C}^r\) 类流形，F 和 \(\mathbf{F}'\) 是以 X 为底空间的 \(\mathbf{C}^r\) 类向量丛，且 \(\phi\) 是从 F 到 \(\mathbf{F}'\) 的态射。对每个 \(x \in \mathbf{X}\)，令 \(\mathbf{S}_x\) 为由下列元素组成的集合：

\[
(a, \phi (a)) \in \mathrm{F} _ {x} \oplus \mathrm{F} _ {x} ^ {\prime}
\]

其中 \(a \in \mathbf{F}_x\)。则 \(\mathbf{S}\) 是各 \(\mathbf{S}_x\) 的并，且是 \(\mathbf{F} \oplus \mathbf{F}'\) 的一个向量子丛。

令 \(\theta\) 和 \(\theta'\) 为从 \(F \oplus F'\) 到自身的映射，定义如下：若 \((u, v) \in F_x \oplus F_x'\)，则

\[
\theta (u, v) = (u, v + \phi (u)), \quad \theta^ {\prime} (u, v) = (u, v - \phi (u)).
\]

由 Differentiable and Analytic Manifolds, R, 7.7.1，\(\theta\) 和 \(\theta'\) 是从 \(F \oplus F'\) 到自身的态射。显然 \(\theta \circ \theta' = \theta' \circ \theta = \mathrm{Id}_{F\oplus F'}\)。因此 \(\theta\) 和 \(\theta'\) 是 \(F \oplus F'\) 的自同构。于是，\(S = \theta(F \oplus \{0\})\) 是 \(F \oplus F'\) 的一个向量子丛。

引理 6. 设 G 是一个李群芽，\(\omega\) 是 G 的左规范微分形式（§ 3, no. 18.9），M 是一个 \(\mathbf{C}^r\) 类流形（\(r \geqslant 2\)），而 \(\alpha\) 是 M 上取值于 L(G) 的 \(\mathbf{C}^{r-1}\) 类 1 次微分形式。

\begin{enumerate}
\def\labelenumi{(\roman{enumi})}
\tightlist
\item
  \(\mathbf{T}(\mathbf{M}\times \mathbf{G})\) 中使微分形式
\end{enumerate}

\[
\theta = \mathrm{pr} _ {1} ^ {*} \alpha - \mathrm{pr} _ {2} ^ {*} \omega
\]

为零的元素构成 T(M × G) 的一个 \(\mathbf{C}^{r-1}\) 类向量子丛 S。

\begin{enumerate}
\def\labelenumi{(\roman{enumi})}
\setcounter{enumi}{1}
\item
  对所有 \((x,g)\in \mathbf{M}\times \mathbf{G},\mathrm{T}(\mathrm{pr}_1)|\mathrm{S}_{(x,g)}\) 都是从 \(\mathrm{S}_{(x,g)}\) 到 \(\mathrm{T}_x(\mathbf{M})\) 的同构。
\item
  若 \(d\alpha + [\alpha]^2 = 0\)（cf.~§3, no. 14），则向量子丛 S 可积。
\end{enumerate}

若 \((x,g)\in \mathbf{M}\times \mathbf{G}\) 且 \((u,v)\in \mathrm{T}_x(\mathrm{M})\times \mathrm{T}_g(\mathrm{G})\)，则

\[
\theta_ {(x, g)} (u, v) = \alpha (u) - g ^ {- 1} v.
\]

因此，\(\theta_{(x,g)}\) 的核是集合 \(\mathrm{S}_{(x,g)}\)，它由 \((u,g\alpha (u))\)（其中 \(u\in \mathrm{T}_x(\mathbf{M})\)）组成，从而得到 (ii)。我们把 \(\mathrm{T}(\mathrm{M}\times \mathrm{G})\) 看作两个向量丛 F 和 \(\mathbf{F}'\) 的直和，其中 \(\mathrm{F}_{(x,g)} = \mathrm{T}_x(\mathrm{M})\times \{0\}\) 且 \(\mathrm{F}_{(x,g)}^{\prime} = \{0\} \times \mathrm{T}_g(\mathrm{G})\)，对所有

\[
(x, g) \in \mathbf {M} \times \mathbf {G}.
\]

对 \(u \in \mathrm{T}_{x}(\mathbf{M}) \times \{0\}\)，记 \(\phi(u) = (0, g\alpha(u))\)。利用 \(\mathrm{T}(\mathrm{G})\) 的左平凡化可见，\(\phi\) 是从 \(\mathbf{F}\) 到 \(\mathbf{F}'\) 的态射，从而得到 (i)（引理 5）。最后，若 \(d\alpha + [\alpha]^2 = 0\)，则

\[
\begin{array}{r l} {d \theta = \mathrm{pr} _ {1} ^ {*} (d \alpha) - \mathrm{pr} _ {2} ^ {*} (d \omega)} \\ & {= - \frac {1}{2} (\mathrm{pr} _ {1} ^ {*} \alpha \wedge \mathrm{pr} _ {1} ^ {*} \alpha - \mathrm{pr} _ {2} ^ {*} \omega \wedge \mathrm{pr} _ {2} ^ {*} \omega)} \\ & {= - \frac {1}{2} (\mathrm{pr} _ {1} ^ {*} \alpha - \mathrm{pr} _ {2} ^ {*} \omega) \wedge (\mathrm{pr} _ {1} ^ {*} \alpha + \mathrm{pr} _ {2} ^ {*} \omega)} \\ & {= - \frac {1}{2} \theta \wedge (\mathrm{pr} _ {1} ^ {*} \alpha + \mathrm{pr} _ {2} ^ {*} \omega)} \end{array}
\]

从而 S 可积（Differentiable and Analytic Manifolds, R, 9.3.6）。

定理 5. 设 \(G\) 是一个李群芽，\(M\) 是一个 \(C^r (r\geqslant 2)\) 类流形，而 \(\alpha\) 是一个 \(\mathbf{C}^{r-1}\) 类 1 次微分形式，定义在 \(M\) 上并取值于 \(L(G)\)，满足 \(d\alpha +[\alpha ]^2 = 0\)。对所有 \(x\in M\) 和 \(g\in G\)，存在一个映射 \(f\)，它属于 \(\mathbf{C}^{r-1}\) 类，在 \(x\) 的某个开邻域上有定义，取值于 \(G\)，满足 \(f(x) = g\) 及 \(f^{-1}.df = \alpha\)。满足这些条件的两个映射在 \(x\) 的某个邻域上重合。

设 \(x \in M\) 且 \(g \in G\)。由引理 6（采用其记号）以及 Differentiable and Analytic Manifolds, R, 9.3.7，存在 M 中 \(x\) 的一个开邻域 U，以及一个 \(m \mapsto \phi(m) = (m, f(m))\) 的 \(C^{r-1}\) 类从 U 到 \(M \times G\) 的映射，使得 \(f(x) = g\) 且 \(\phi^*(\theta) = 0\)。于是

\[
\begin{array}{l l} f ^ {- 1}. d f = f ^ {*} (\omega) & \text {(§ 3, no. 18.9)} \\ = (\mathrm{pr} _ {2} \circ \phi) ^ {*} (\omega) & \text {(for \(f = \mathrm{pr} _ {2} \circ \phi\))} \\ = \phi^ {*} (\mathrm{pr} _ {1} ^ {*} \alpha - \theta) & \text {(Lemma 6)} \\ = \phi^ {*} (\mathrm{pr} _ {1} ^ {*} \alpha) & \text {(for \(\phi^ {*} (\theta) = 0\))} \\ = \alpha & \text {(for \(\mathrm{pr} _ {1} \circ \phi = \mathrm{Id} _ {\mathbf {U}}\))}. \end{array}
\]

设 \(f'\) 是从 U 到 G 的 \(\mathbf{C}^{r-1}\) 类映射，满足 \(f'(x) = g\) 和 \(f'^{-1}df' = \alpha\)。由 § 3, 18.9，\(ff'^{-1}\) 局部为常值，因而 \(f' = f\) 在 \(x\) 的某个邻域上成立。

命题 10. 设 M 是解析流形，g 是完备可赋范李代数，\(\alpha\) 是 M 上取值于 g 的解析 1 次微分形式，并具有以下性质：

\begin{enumerate}
\def\labelenumi{(\alph{enumi})}
\item
  对所有 \(m \in \mathbf{M}\)，\(\alpha_{m}\) 都是从 \(\mathrm{T}_{m}(\mathrm{M})\) 到 \(\mathfrak{g}\) 的同构；
\item
  \(d\alpha + [\alpha]^{2} = 0.\)
\end{enumerate}

那么，对所有 \(m_0 \in \mathbf{M}\)，存在一个开邻域 \(\mathbf{M}'\)，它是 \(m_0\) 在 \(\mathbf{M}\) 中的邻域，并且 \(\mathbf{M}'\) 上带有与 \(\mathbf{M}'\) 的流形结构相容的李群芽结构，单位元为 \(m_0\)，并具有以下性质：

\begin{enumerate}
\def\labelenumi{(\roman{enumi})}
\item
  \(\alpha_{m_{0}}\) 是从 L(M') 到 g 的同构；
\item
  微分形式 \(m \mapsto \alpha_{m_0}^{-1} \circ \alpha_m\) 是 \(M'\) 的左规范微分形式。若 \(M_1'\) 和 \(M_2'\) 是两个这样的群芽，则 \(M_1'\) 和 \(M_2'\) 有一个公共开子群芽。
\end{enumerate}

存在一个李群芽 G，使得 \(L(G) = g\)。令 \(m_{0} \in M\)。由定理 5，存在一个开邻域 \(M'\)（\(m_{0}\) 在 M 中），以及一个从 \(M'\) 到 G 的解析映射 f，使得 \(f(m_{0}) = e\) 且 \(f^{-1}.df = \alpha\)。于是 \(T_{m_{0}}(f) = \alpha_{m_{0}}\) 是从 \(T_{m_{0}}(M)\) 到 g 的同构；因此，缩小 \(M'\) 和 G 后，可以设 f 是从流形 \(M'\) 到流形 G 的同构。把 G 上的李群芽结构传递给 \(M'\)，所用的同构为 \(f^{-1}\)。于是 \(T_{m_{0}}(f)\) 成为从 \(L(M')\) 到 \(L(G) = g\) 的同构，从而得到 (i)。另一方面，若 \(\omega\) 表示 G 的左规范微分形式，则

\[
\begin{array}{r} \alpha_ {m _ {0}} ^ {- 1} \circ \alpha_ {m} = (\mathrm{T} _ {m _ {0}} f) ^ {- 1} \circ (f ^ {- 1}. d f) (m) \\ = (\mathrm{T} _ {m} f) ^ {- 1} \circ \omega (f (m)) \circ \mathrm{T} _ {m} f \end{array}
\]

从而 \(m \mapsto \alpha_{m_0}^{-1} \circ \alpha_m\) 是 \(\mathbf{M}'\) 的左规范微分形式。

令 \(\mathbf{M}^{\prime}\) 为 \(m_0\) 的一个开邻域，其上带有李群芽结构，单位元为 \(m_0\)，并具有与性质 (i) 和 (ii) 类似的性质。则 \(\alpha_{m_0}\) 既是从 \(\mathrm{L}(\mathbf{M}')\) 到 \(\mathfrak{g}\) 的同构，也是从 \(\mathrm{L}(\mathbf{M}'')\) 到 \(\mathfrak{g}\) 的同构，因而 \(\mathrm{L}(\mathbf{M}') = \mathrm{L}(\mathbf{M}'')\)。因此，缩小 \(\mathbf{M}'\) 和 \(\mathbf{M}''\) 后，可以设存在一个同构 \(\phi\)，它把群芽 \(\mathbf{M}'\) 映到群芽 \(\mathbf{M}''\)（no. 1, Corollary 1 to Theorem 1）。于是 \(\phi^{-1}.d\phi\) 是 \(\mathbf{M}'\) 的左规范微分。另一方面，令 \(\psi\) 为从流形 \(\mathbf{M}' \cap \mathbf{M}''\) 到李群芽 \(\mathbf{M}''\) 的规范嵌入；显然 \(\psi^{-1}.d\psi\) 是 \(\mathbf{M}''\) 的左规范微分的限制。因此有 \((\psi^{-1}.d\psi)(m) = \alpha_{m_0}^{-1} \circ \alpha_m = (\phi^{-1}.d\phi)(m)\)，对所有 \(m \in \mathbf{M}' \cap \mathbf{M}''\) 成立。所以 \(\phi\) 和 \(\psi\) 在 \(m_0\) 的某个邻域上重合（§ 3, 18.9）。这就证明了命题的最后一个断言。

推论。设 M 是有限维解析流形。~令 \(\omega_{1},\ldots,\omega_{n}\) 为 M 上取值为标量的解析 1 次微分形式，它们在 M 的每一点都线性无关，并且对所有 \(k=1,\ldots,n\)，\(\omega_{k}\) 都是 \(\omega_{i}\wedge\omega_{j}\) 的常系数线性组合。那么，对所有 \(m_{0}\in M\)，存在一个开邻域 \(M'\)，它是 \(m_{0}\) 在 M 中的邻域，并且带有定义在 \(M'\) 上、与 \(M'\) 的流形结构相容的李群芽结构，单位元为 \(m_{0}\)，使得 \(\omega_{1}|_{M'},\ldots,\omega_{n}|_{M'}\) 构成取值为标量的左不变 1 次微分形式空间在 \(M'\) 上的一组基。

若 \(\mathbf{M}_1'\) 和 \(\mathbf{M}_2'\) 是两个这样的群芽，则 \(\mathbf{M}_1'\) 和 \(\mathbf{M}_2'\) 有一个公共开子群芽。

令 \(\mathbf{X}_1, \ldots, \mathbf{X}_n\) 为 M 上的向量场，使得在 M 的每一点 \(m\)，\((\mathbf{X}_i)_m\) 构成 \(\mathrm{T}_m(\mathrm{M})\) 的一组与 \(((\omega_1)_m), \ldots, (\omega_n)_m)\) 对偶的基。这些向量场是解析的。由假设，存在 \(c_{ijk} \in \mathbf{K}\)（\(1 \leqslant i,j, k \leqslant n\)），使 \(c_{ijk} = -c_{jik}\) 且 \(d\omega_k = \sum_{i < j} c_{ijk}\omega_i \wedge \omega_j\)。由 Differentiable and Analytic Manifolds, R, 8.5.7，公式 (11)，

\[
\langle [ \mathrm{X} _ {i}, \mathrm{X} _ {j} ], \omega_ {k} \rangle = - (d \omega_ {k}) (\mathrm{X} _ {i}, \mathrm{X} _ {j}) = - \left(\sum_ {r <   s} c _ {r s k} \omega_ {r} \wedge \omega_ {s}\right) (\mathrm{X} _ {i}, \mathrm{X} _ {j}) = - c _ {i j k}
\]

从而 \([\mathrm{X}_i,\mathrm{X}_j] = -\sum_{k}c_{ijk}\mathrm{X}_k\)。于是 \(-c_{ijk}\) 是李代数 \(\mathfrak{g}\) 相对于基 \((e_1,\dots ,e_n)\) 的结构常数。对所有 \(m\in \mathbf{M}\)，令 \(\alpha_{m}\) 为从 \(\mathrm{T}_{m}(\mathrm{M})\) 到 \(\mathfrak{g}\) 的线性映射，它把 \((\mathrm{X}_1)_m\) 映到 \(e_1,\ldots ,(\mathrm{X}_n)_m\) 映到 \(e_n\)。于是 \(\alpha\) 是 \(\mathbf{M}\) 上的解析微分形式，次数为

1，取值于 \(\mathfrak{g}\)，且 \(\alpha_{m}\) 是从 \(\mathrm{T}_m(\mathbf{M})\) 到 \(\mathfrak{g}\) 的同构。另一方面，\(\alpha = \sum_{k=1}^{n} \omega_k e_k\)，因此

\[
d \alpha = \sum_ {k = 1} ^ {n} (d \omega_ {k}) e _ {k} = \sum_ {k = 1} ^ {n} \left(\sum_ {i <   j} c _ {i j k} \omega_ {i} \wedge \omega_ {j}\right) e _ {k}
\]

并且

\[
\begin{array}{r l} {[ \alpha ] ^ {2} = \sum_ {k = 1} ^ {n} [ \omega_ {k} e _ {k} ] ^ {2} + \sum_ {i <   j} (\omega_ {i} e _ {i}) \wedge (\omega_ {j} e _ {j})} & {\text {(§ 3, formula (30))}} \\ {= \sum_ {i <   j} (\omega_ {i} \wedge \omega _ {j}) [ e _ {i}, e _ {j} ]} \\ {= - \sum_ {k = 1} ^ {n} \sum_ {i <   j} (c _ {i j k} \omega_ {i} \wedge \omega_ {j}) e _ {k}} \\ {= - d \alpha.} \end{array}\tag{5}
\]

只需应用命题 10。

\subsection{从无穷小作用律到作用律的过渡}

命题 11. 设 \(\mathrm{G}_1\) 和 \(\mathrm{G}_2\) 是李群芽，\(\mathbf{X}_1\) 和 \(\mathbf{X}_2\) 是 \(\mathrm{C}^r\) 类流形（\(r \geqslant 2\)）。对于 \(i = 1, 2\)，令 \(\psi_i\) 为 \(\mathrm{C}^r\) 类 \(\mathrm{G}_i\) 在 \(\mathbf{X}_i\) 上的左作用律片段，\(\mathrm{D}_i\) 为相应的无穷小作用律。令 \(\mu: \mathrm{G}_1 \to \mathrm{G}_2\) 为态射，\(\phi: \mathrm{X}_1 \to \mathrm{X}_2\) 为 \(\mathrm{C}^r\) 类映射。假设对所有 \(a \in \mathrm{L}(\mathrm{G})\)，向量场 \((\mathrm{D}_{1})_a\) 和 \((\mathrm{D}_{2})_{\mathrm{L}(\mu)a}\) 关于 \(\phi\) 相关（Differentiable and Analytic Manifolds, R, 8.2.6）。则存在一个邻域 \(\Omega\)，它包含 \(\{e\} \times \mathbf{X}_1\) 且位于 \(\mathrm{G}_1 \times \mathbf{X}_1\) 中，使得 \(\phi(\psi_1(g,x)) = \psi_2(\mu(g),\phi(x))\) 对所有 \((g,x) \in \Omega\) 成立。

令 \(p_1: \mathrm{G}_1 \times \mathrm{X}_2 \to \mathrm{G}_1\)、\(p_2: \mathrm{G}_1 \times \mathrm{X}_2 \to \mathrm{X}_2\) 为规范投影。对所有 \((g_1, x_2) \in \mathrm{G}_1 \times \mathrm{X}_2\)，令 \(f_{g_1, x_2}\) 为映射 \(g_1 a \mapsto (\mathrm{D}_2)_{\mathrm{L}(\mu)a}(x_2)\)，它从 \(\mathrm{T}_{g_1}(\mathrm{G}_1) = g_1\mathrm{L}(\mathrm{G}_1)\) 到 \(\mathrm{T}_{x_2}(\mathrm{X}_2)\)。这些 \(f_{g_1, x_2}\) 定义了从 \(p_1^* \mathrm{T}(\mathrm{G}_1)\) 到 \(p_2^* \mathrm{T}(\mathbf{X}_2)\) 的态射。

令 \(x_0 \in X_1\)。存在 e 的一个开邻域 G（位于 \(G_1\) 中）以及 \(x_0\) 在 \(X_1\) 中的一个开邻域 X，使得 \(\psi_1(g, x)\) 和 \(\psi_2(\mu(g), \phi(x))\) 对 \((g, x) \in G \times X\) 有定义。对于 \((g, x) \in G \times X\)，记

\[
\alpha (g, x) = \phi (\psi_ {1} (g, x)) \in \mathrm{X} _ {2}, \quad \beta (g, x) = \psi_ {2} (\mu (g), \phi (x)) \in \mathrm{X} _ {2}.
\]

若 \(\mathbf{G}\) 和 \(\mathbf{X}\) 足够小，则对所有 \((a,g,x)\in \mathrm{L}(\mathrm{G}_1)\times \mathrm{G}\times \mathrm{X},\)

\[
\begin{array}{r l} (\mathrm{T} \alpha) (a g, 0 _ {x}) & = (\mathrm{T} \phi) ((\mathrm{D} _ {1}) _ {a} (\psi_ {1} (g, x))) \\ & = (\mathrm{D} _ {2}) _ {L (\mu) a} (\phi (\psi_ {1} (g, x))) \\ & = (\mathrm{D} _ {2}) _ {L (\mu) a} \alpha (g, x), \end{array}
\]

\[
\begin{array}{r l} (\mathrm{T} \beta) (a g, 0 _ {x}) = & (\mathrm{T} \psi_ {2}) (\mathrm{L} (\mu) a. \mu (g), \phi (x)) \\ = & (\mathrm{D} _ {2}) _ {\mathrm{L} (\mu) a} (\psi_ {2} (\mu (g), \phi (x))) \\ = & (\mathrm{D} _ {2}) _ {\mathrm{L} (\mu) a} \beta (g, x). \end{array}
\]

因此对 \(x \in \mathbf{X}\)，映射 \(g \mapsto \alpha(g, x)\) 和 \(g \mapsto \beta(g, x)\) 都是 \(f\) 的积分；由于

\[
\beta (e, x) = \phi (x) = \alpha (e, x)
\]

对所有 \(x \in \mathbf{X}\)，由 Differentiable and Analytic Manifolds, R, 9.3.7，可知 \(\alpha\) 和 \(\beta\) 在 \((e, x_0)\) 的某个邻域上重合。因此命题得证。

推论。设 G 是李群芽，X 是 \(\mathbf{C}^{r}\) 类流形。考虑 G 在 X 上的两个 \(\mathbf{C}^{r}\) 类左作用律片段。假设对所有 \(a \in L(G)\)，相应的 X 上的向量场 \(D_{a}\) 对两个作用律片段都相同。则这两个作用律片段在 \(\{e\} \times X\) 的某个邻域上重合。

定理 6. 设 G 是李群芽，X 是 \(\mathbf{C}^r\) 类流形（\(r \geqslant 2\)），\(x_0\) 是 X 中的一点。令 \(a \mapsto \mathrm{D}_a\) 为一个 \(\mathbf{C}^{r-1}\) 类的、在 X 上的 \(\mathrm{L(G)}\) 左无穷小作用律。

\begin{enumerate}
\def\labelenumi{(\roman{enumi})}
\item
  存在一个开邻域 \(\mathbf{X}'\)，它包含 \(x_0\) 且位于 \(\mathbf{X}\) 中，并且存在一个 \(\mathbf{C}^{r-1}\) 类的、由 \(\mathbf{G}\) 在 \(\mathbf{X}'\) 上定义的左作用律片段，使其相应的无穷小作用律为 \(a \mapsto \mathbf{D}_a|\mathbf{X}'\)。
\item
 设有两个 \(\mathbf{C}^{r-1}\) 类的、由 \(\mathbf{G}\) 在 \(\mathbf{X}''\)（它是 \(x_0\) 的一个开邻域）上定义的左作用律片段；若它们都以 \(a \mapsto \mathbf{D}_a|\mathbf{X}''\) 为相应的无穷小作用律，则它们在 \((e, x_0)\) 的某个邻域上重合。
\end{enumerate}

断言 (ii) 由命题 11 的推论推出。下面证明 (i)。对所有 \((g, x) \in \mathbf{G} \times \mathbf{X}\) 和 \(a \in \mathrm{L}(\mathbf{G})\)，记

\[
\mathrm{Q} _ {a} (g, x) = (a g, \mathrm{D} _ {a} (x)) \in \mathrm{T} _ {g} (\mathrm{G}) \times \mathrm{T} _ {x} (\mathrm{X}).
\]

令 \(S_{(g,x)}\) 为由 \(Q_{a}(g,x)\)（其中 \(a\in L(G)\)）组成的集合。由第 6 号的引理 5，\(S_{(g,x)}\) 是 \(T(G)\times T(X)\) 的一个向量子丛 S 的纤维。令 a, b 属于 \(L(G)\)；则

\[
\begin{array}{r l} {[ \mathrm{Q} _ {a}, \mathrm{Q} _ {b} ] (g, x) = ([ \mathrm{R} _ {a}, \mathrm{R} _ {b} ] (g), [ \mathrm{D} _ {a}, \mathrm{D} _ {b} ] (x))} \\ & = (- \mathrm{R} _ {[ a, b ]} (g), - \mathrm{D} _ {[ a, b ]} (x)) \\ & = \mathrm{Q} _ {- [ a, b ]} (g, x) \end{array}\tag{§ 3, 18.6}
\]

从而 S 可积（Differentiable and Analytic Manifolds, R, 9.3.3, (iv)）。由 Differentiable and Analytic Manifolds, R, 9.3.7，存在一个开邻域 \(\mathrm{G}_1\)，它是 G 中 \(e\) 的邻域；一个开邻域 \(\mathbf{X}_1\)，它是 X 中 \(x_0\) 的邻域；以及一个映射 \((g,x)\mapsto gx\)，它是 \(\mathrm{C}^{r - 1}\) 类，从 \(\mathrm{G}_1\times \mathrm{X}_1\) 映到 X，使得 \(ex = x\) 对所有 \(x\in \mathbf{X}_1\) 都成立，并且

\[
(a g) x = \mathrm{D} _ {a} (g x) \quad \text { for } a \in \mathrm{L} (\mathrm{G}), g \in \mathrm{G} _ {1}, x \in \mathrm{X} _ {1}.\tag{6}
\]

特别地

\[
a x = \mathrm{D} _ {a} (x).\tag{7}
\]

令 \(G_{2}\) 为 \(G_{1}\) 中 e 的一个开邻域，\(X_{2}\) 为一个开邻域，它包含 \(x_{0}\) 且位于 \(X_{1}\) 中，使得 \(gg'\) 有定义且属于 \(G_{1}\)，其中 g、\(g'\) 属于 \(G_{2}\)，并且 gx 有定义且属于 \(\mathbf{X}_1\)，其中 \((g,x)\in \mathbf{G}_2\times \mathbf{X}_2\)。考虑映射 \(\alpha_{1},\alpha_{2}\)，它们从 \(\mathrm{G}_2\times (\mathrm{G}_2\times \mathrm{X}_2)\) 到 X，定义为

\[
\alpha_ {1} (g, (h, x)) = g (h x), \quad \alpha_ {2} (g, (h, x)) = (g h) x.
\]

它们属于 \(\mathbf{C}^{r - 1}\) 类。于是

\[
\alpha_ {1} (e, (h, x)) = h x = \alpha_ {2} (e, (h, x)).
\]

另一方面

\[
\begin{array}{r l} \mathrm{T} (\alpha_ {1}) (a g, 0 _ {(h, x)}) & = (a g) (h x) \\ & = \mathrm{D} _ {a} (g (h x)) \qquad \text {by (6)} \\ & = \mathrm{D} _ {a} (\alpha_ {1} (g, (h, x))), \\ \mathrm{T} (\alpha_ {2}) (a g, 0 _ {(h, x)}) & = (a g h) x \\ & = \mathrm{D} _ {a} ((g h) x) \qquad \text {by (6)} \\ & = \mathrm{D} _ {a} (\alpha_ {2} (g, (h, x))). \end{array}
\]

由 Differentiable and Analytic Manifolds, R, 9.3.7，\(\alpha_{1}\) 和 \(\alpha_{2}\) 在 \((e, (e, x_{0}))\) 的某个邻域上重合。于是 (i) 由 § 1, no. 11 的 Proposition 23 和 (7) 得出。

推论 1. 设 G 是李群芽，X 是 \(\mathbf{C}^r\) 类旁紧流形（\(r \geqslant 2\)）。令 \(a \mapsto \mathrm{D}_a\) 为 L(G) 在 X 上的 \(\mathbf{C}^{r-1}\) 类左无穷小作用律。

\begin{enumerate}
\def\labelenumi{(\roman{enumi})}
\item
  存在一个 \(\mathbf{C}^{r-1}\) 类的、由 \(\mathbf{G}\) 在 \(\mathbf{X}\) 上定义的左作用律片段，使其相应的无穷小作用律为 \(a \mapsto \mathbf{D}_a\)。
\item
  两个 \(\mathbf{C}^{r-1}\) 类的 G 在 X 上的左作用律，若都以 \(a \mapsto \mathbf{D}_a\) 为相应的无穷小作用律，则它们在 \(\{e\} \times \mathbf{X}\) 的某个邻域上重合。
\end{enumerate}

断言 (ii) 由命题 11 的推论推出。由定理 6 (i)，存在一个开覆盖 \((\mathbf{X}_i)_{i\in I}\)，它覆盖 \(\mathbf{X}\)，并且对所有 \(i\in I\)，有一个左作用律片段 \(\psi_i\)，它是 \(\mathbf{C}^{r-1}\) 类的，由 \(\mathbf{G}\) 定义在 \(\mathbf{X}_i\) 上，使相应的无穷小作用律为 \(a\mapsto \mathrm{D}_a|\mathbf{X}_i\)。由于 \(\mathbf{X}\) 是旁紧的，可以设覆盖 \((\mathbf{X}_i)_{i\in I}\) 是局部有限的。对所有 \((i,j)\in I\times I\) 及所有 \(x\in \mathbf{X}_i\cap \mathbf{X}_j\)，\(\psi_i\) 和 \(\psi_j\) 在 \((e,x)\) 的某个邻域上重合（命题 11 的推论）。由于 \(\mathbf{X}\) 是正规空间，我们可以应用 § 1, no. 11 的 Proposition 24，从而证明 (i)。

推论 2. 设 X 是 \(\mathbf{C}^r\) 类旁紧流形（\(r \geqslant 2\)），\(\xi\) 是 X 上的 \(\mathbf{C}^{r-1}\) 类向量场。存在 K 在 X 上的作用律片段 \(\psi\)，它是 \(\mathbf{C}^{r-1}\) 类的，使得对所有 \(x \in \mathbf{X}\)，\(\xi(x)\) 是 K 在 0 处的切向量 1 在 \(t \mapsto (t, x)\) 下的像。具有上述性质的两个作用律片段在 \(\{0\} \times \mathbf{X}\) 的某个邻域上重合。

这是推论 1 的特殊情形。

注记。在本号中，左作用律当然可以处处替换为右作用律。

\section{李群中的形式计算}

令 \(f, g\) 为关于同一组不定元、系数属于 K 的两个形式幂级数，令 \(f_i\)（分别地，\(g_i\)）为 \(f\)（分别地，\(g\)）的 i 次齐次分量。记

\[
f \equiv g \mod \deg p
\]

\[
\text { if } f _ {i} = g _ {i} \text { for } i <   p.
\]

在本节中，G 表示一个 n 维李群芽；设基域 K 的特征为零。我们一次性借助一个坐标图，把 G 中 e 的一个开邻域同 \(K^{n}\) 中 0 的一个开邻域 U 识别，从而 e 识别为 0。对于 U 中的 x, y 以及 \(n \in \mathbf{Z}\)，x.y 表示 x 与 y 的乘积，\(x^{[m]}\) 表示 G 中 x 的 m 次幂（在它们有定义时）。\(x \in U\) 的坐标记为 \(x_{1}, x_{2}, \ldots, x_{n}\)。

\subsection{\texorpdfstring{系数 \(c_{\alpha\beta\gamma}\)}{系数 c\_\{\textbackslash alpha\textbackslash beta\textbackslash gamma\}}}

令 \(\Omega\) 为所有 \((x,y)\in\mathbf{U}\times\mathbf{U}\) 的集合，其中 \(x.y\) 有定义且属于 \(\mathbf{U}\)。则 \(\Omega\) 是 \(\mathbf{U}\times\mathbf{U}\) 中的开集，且映射 \((x,y)\mapsto x.y\) 从 \(\Omega\) 到 \(\mathbf{U}\) 是解析的。因此，\(z_{1},\ldots,z_{n}\)（即 \(z=x.y\) 的坐标）可在原点附近按 \(x_{1},\ldots,x_{n},y_{1},\ldots,y_{n}\) 的幂展开为幂级数。因此，存在定义良好的常数 \(c_{\alpha_{1},\ldots,\alpha_{n},\beta_{1},\ldots,\beta_{n}}\in\mathbf{K}\)，使得

\[
z _ {1} ^ {\gamma_ {1}} \dots z _ {n} ^ {\gamma_ {n}} = \sum_ {\alpha_ {1},\dots,\alpha_ {n},\beta_ {1},\dots,\beta_ {n} \in \mathbf {N}} c _ {\alpha_ {1},\dots,\alpha_ {n},\beta_ {1},\dots,\beta_ {n},\gamma_ {1},\dots,\gamma_ {n}} x _ {1} ^ {\alpha_ {1}} \dots x _ {n} ^ {\alpha_ {n}} y _ {1} ^ {\beta_ {1}} \dots y _ {n} ^ {\beta_ {n}}\tag{1}
\]

对于 \(\gamma_{1},\ldots,\gamma_{n}\) 属于 \(\mathbf{N}\)。采用 Differentiable and Analytic Manifolds, R 中的约定，我们把这些公式简写为：

\[
(x. y) ^ {\gamma} = \sum_ {\alpha , \beta \in \mathbf {N} ^ {n}} c _ {\alpha , \beta , \gamma} x ^ {\alpha} y ^ {\beta} \quad (\gamma \in \mathbf {N} ^ {n}).\tag{2}
\]

由于 \(x.0 = 0.x = x\) 对 \(x\in \mathbf{U}\) 成立，

\[
c _ {\alpha , 0, \gamma} = c _ {0, \alpha , \gamma} = \delta_ {\alpha \gamma}\tag{3}
\]

其中 \(\delta_{\alpha \gamma}\) 是克罗内克指标。特别地，以下从现在起用 \(k\) 代替 \(\varepsilon_k\)，其中 \(k = 1, \ldots, n\)，

\[
(x. y) _ {k} = x _ {k} + y _ {k} + \sum_ {| \alpha | \geqslant 1, | \beta | \geqslant 1} c _ {\alpha , \beta , k} x ^ {\alpha} y ^ {\beta}.\tag{4}
\]

记 \(c_{\alpha \beta} = (c_{\alpha \beta 1}, c_{\alpha \beta 2}, \ldots, c_{\alpha \beta n}) \in \mathbf{K}^n\)，于是得到

\[
x. y = x + y + \sum_ {| \alpha | \geqslant 1, | \beta | \geqslant 1} c _ {\alpha \beta} x ^ {\alpha} y ^ {\beta}.\tag{5}
\]

在 (5) 的右端，我们取 2 次齐次分量：

\[
\mathrm{B} (x, y) = \sum_ {i, j = 1} ^ {n} c _ {i j} x _ {i} y _ {j}
\]

于是 \((x,y)\mapsto \mathrm{B}(x,y)\) 是从 \(\mathbf{K}^n\times \mathbf{K}^n\) 到 \(\mathbf{K}^n\) 的双线性映射。于是

\[
x. y \equiv x + y + \mathrm{B} (x, y) \mod \deg 3.\tag{6}
\]

另一方面，公式 (4) 蕴含

\[
c _ {\alpha , \beta , \gamma} = 0 \quad \text { if } | \alpha | + | \beta | <   | \gamma |\tag{7}
\]

并且展开式中总次数为 \(|\gamma|\) 的、\(z^\gamma\) 的项也正是 \((x_1 + y_1)^{\gamma_1}(x_2 + y_2)^{\gamma_2} \ldots (x_n + y_n)^{\gamma_n} = \sum_{\alpha + \beta = \gamma} ((\alpha, \beta)) x^\alpha y^\beta\) 的相应项（cf.~Differentiable and Analytic Manifolds, R, Notation and conventions）。因此：

\begin{enumerate}
\def\labelenumi{(\arabic{enumi})}
\setcounter{enumi}{7}
\tightlist
\item
\end{enumerate}

\[
c _ {\alpha , \beta , \gamma} = 0 \quad \text { if } \quad | \alpha | + | \beta | = | \gamma | \quad \text { but } \quad \alpha + \beta \neq \gamma\tag{8}
\]

\[
c _ {\alpha , \beta , \alpha + \beta} = ((\alpha , \beta)).\tag{9}
\]

乘法的结合律蕴含关系

\[
\sum_ {\alpha , \xi} c _ {\alpha \xi \eta} x ^ {\alpha} \left(\sum_ {\beta , \gamma} c _ {\beta \gamma \xi} y ^ {\beta} z ^ {\gamma}\right) = \sum_ {\xi , \gamma} c _ {\xi \gamma \eta} \left(\sum_ {\alpha , \beta} c _ {\alpha \beta \xi} x ^ {\alpha} y ^ {\beta}\right) z ^ {\gamma}
\]

对充分接近 0 的所有 \(x, y, z\)，从而

\[
\sum_ {\xi} c _ {\alpha \xi \eta} c _ {\beta \gamma \xi} = \sum_ {\xi} c _ {\xi \gamma \eta} c _ {\alpha \beta \xi} (\alpha , \beta , \gamma , \eta \text { in } \mathbf {N} ^ {n}).\tag{10}
\]

李群芽 G 有交换开子群芽，当且仅当 \(c_{\alpha\beta\gamma} = c_{\beta\alpha\gamma}\) 对所有 \(\alpha, \beta, \gamma\)（它们属于 \(\mathbf{N}^{n}\)）都成立。

\subsection{李代数中的括号}

对于 \(\alpha \in \mathbf{N}^n\)，令 \(e_{\alpha}\) 为原点处的点分布 \(\frac{1}{\alpha!} \frac{\partial^{\alpha}}{\partial x^{\alpha}}\)。特别地，\(e_k = e_{\varepsilon_k} = \frac{\partial}{\partial x_k}\)。这些 \(e_{\alpha}\) 构成向量空间 \(\mathrm{U}(G)\) 的一组基。若 \(f\) 是 G 中 0 的一个开邻域上的解析函数，且 \(f(x) = \sum_{\alpha} \lambda_{\alpha} x^{\alpha}\) 是它在原点附近的幂级数展开，则

特别地，

\[
\langle e _ {\alpha}, f \rangle = \lambda_ {\alpha}.
\]

\[
\langle e _ {\alpha}, x ^ {\gamma} \rangle = \delta_ {\alpha \gamma}.
\]

因此

\[
\begin{array}{r l} \langle e _ {\alpha} * e _ {\beta}, x ^ {\gamma} \rangle & = \langle e _ {\alpha} \otimes e _ {\beta}, (x. y) ^ {\gamma} \rangle \\ & = \langle e _ {\alpha} \otimes e _ {\beta}, \sum_ {\alpha^ {\prime}, \beta^ {\prime}} c _ {\alpha^ {\prime} \beta^ {\prime} \gamma} x ^ {\alpha^ {\prime}} y ^ {\beta^ {\prime}} \rangle \\ & = \sum_ {\alpha^ {\prime}, \beta^ {\prime}} c _ {\alpha^ {\prime} \beta^ {\prime} \gamma} \langle e _ {\alpha}, x ^ {\alpha^ {\prime}} \rangle \langle e _ {\beta}, y ^ {\beta^ {\prime}} \rangle = c _ {\alpha \beta \gamma} \end{array}
\]

从而

\[
e _ {\alpha} * e _ {\beta} = \sum_ {\gamma} c _ {\alpha \beta \gamma} e _ {\gamma}.\tag{11}
\]

（公式 (10) 表示 U(G) 上的结合律。）

特别地，由于 L(G) 在括号运算下稳定，

\[
[ e _ {i}, e _ {j} ] = \sum_ {k} (c _ {i j k} - c _ {j i k}) e _ {k}.\tag{12}
\]

因此，\(\mathbf{L}(\mathbf{G})\) 相对于基 \((e_{1},\ldots,e_{n})\) 的结构常数为 \(c_{ijk}-c_{jik}\)。换言之，将 \(\mathbf{L}(\mathbf{G})\) 规范地识别为 \(K^{n}\)，得到：

\[
[ x, y ] = \mathrm{B} (x, y) - \mathrm{B} (y, x).\tag{13}
\]

命题 1.

\begin{enumerate}
\def\labelenumi{(\roman{enumi})}
\tightlist
\item
\end{enumerate}

\[
x ^ {[ - 1 ]} \equiv - x + \mathrm{B} (x, x) \mod \deg 3\tag{i}
\]

\[
x. y. x ^ {[ - 1 ]} \equiv y + [ x, y ] \quad \text {   mod   } \deg 3\tag{ii}
\]

\[
y ^ {[ - 1 ]}. x. y \equiv x + [ x, y ] \quad \text {   mod   } \deg 3\tag{iii}
\]

\[
x ^ {[ - 1 ]}. y ^ {[ - 1 ]}. x. y \equiv [ x, y ] \quad \text {   mod   } \deg 3\tag{iv}
\]

\[
x. y. x ^ {[ - 1 ]}. y ^ {[ - 1 ]} \equiv [ x, y ] \quad \text {   mod   } \deg 3\tag{v}.
\]

（在 (i) 中，\(x^{[-1]}\) 表示函数 \(x \mapsto x^{[-1]}\) 在原点附近的幂级数展开；其他公式可以类似解释。）

令 \(g_{1}\) 和 \(g_{2}\) 为 \(x^{[-1]}\) 的 1 次和 2 次齐次分量。于是

\[
\begin{array}{r l} 0 = x. x ^ {[ - 1 ]} \\ \equiv x + g _ {1} (x) + \mathrm{B} (x, g _ {1} (x)) & \text {   mod   } \deg 2 \quad \text {(by (6))} \\ \equiv x + g _ {1} (x) & \text {   mod   } \deg 2 \end{array}
\]

从而 \(g_{1}(x) = -x\)。于是

\[
\begin{array}{r l} 0 & = x. x ^ {[ - 1 ]} \\ & \equiv x + (- x + g _ {2} (x)) + B (x, - x + g _ {2} (x)) \mod \deg 3 \\ & \equiv g _ {2} (x) - B (x, x) \mod \deg 3 \end{array}
\]

从而 \(g_{2}(x) = \mathrm{B}(x, x)\)。这证明了 (i)。然后利用 (i)，

\[
\begin{array}{l l} x. y. x ^ {[ - 1 ]} \equiv (x + y + \mathrm{B} (x, y)) \cdot (- x + \mathrm{B} (x, x)) & \text {mod} \deg 3 \\ \equiv x + y + \mathrm{B} (x, y) + (- x + \mathrm{B} (x, x)) + \mathrm{B} (x + y, - x) & \text {mod} \deg 3 \\ \equiv y + \mathrm{B} (x, y) - \mathrm{B} (y, x) & \text {mod} \deg 3 \\ \equiv y + [ x, y ] \qquad \qquad \qquad \qquad \qquad \qquad \qquad \qquad \qquad \qquad \qquad \qquad \qquad \qquad \qquad \qquad \qquad \qquad \qquad \qquad \qquad \qquad \qquad \qquad \qquad \qquad \qquad \qquad \qquad \qquad \qquad \qquad \qquad \qquad \qquad \qquad \qquad \\ & (\text {by (13)}) \end{array}
\]

从而得到 (ii)。(iii) 的证明类似。结合 (i) 和 (iii)，得到

\[
\begin{array}{r l r} x ^ {[ - 1 ]}. y ^ {[ - 1 ]}. x. y & \equiv (- x + \mathrm{B} (x, x)). (x + [ x, y ]) & \bmod \deg 3 \\ & \equiv - x + \mathrm{B} (x, x) + x + [ x, y ] + \mathrm{B} (- x, x) & \bmod \deg 3 \\ & \equiv [ x, y ] & \bmod \deg 3 \end{array}
\]

从而得到 (iv)。(v) 的证明类似。

\subsection{幂}

考虑 \(j\) 个 \(\mathbf{G}\) 中的点：

\[
\begin{array}{l} x (1) = (x (1) _ {1}, x (1) _ {2}, \ldots , x (1) _ {n}) \\ x (2) = (x (2) _ {1}, x (2) _ {2}, \ldots , x (2) _ {n}) \\ \cdot \cdot \cdot \\ x (j) = (x (j) _ {1}, x (j) _ {2}, \ldots , x (j) _ {n}). \end{array}
\]

映射 \((x(1), x(2), \ldots, x(j)) \mapsto x(1).x(2) \ldots x(j)\) 在原点附近可展开为幂级数：

\[
x (1). x (2) \dots x (j) = \sum_ {\alpha (1), \alpha (2), \dots , \alpha (j) \in \mathbf {N} ^ {n}} a _ {\alpha (1), \dots , \alpha (j)} x (1) ^ {\alpha (1)} \dots x (j) ^ {\alpha (j)}\tag{14}
\]

其中 \(a_{\alpha(1), \ldots, \alpha(j)}\) 是 \(\mathbf{K}^n\) 的元素。对于 \(j = 0, 1, 2, \ldots\)，记

\[
\psi_ {j} (x) = \sum_ {\alpha (1) \neq 0, \dots , \alpha (j) \neq 0} a _ {\alpha (1), \dots , \alpha (j)} x ^ {\alpha (1) + \dots + \alpha (j)}\tag{15}
\]

其中右端是关于变量 \(x \in K^{n}\) 的收敛幂级数。这个级数由 (14) 中删去某个变量 \(x(1), \ldots, x(j)\) 未显式出现的项，然后令 \(x(1) = x(2) = \cdots = x(j) = x\) 而得到。

若 \(t \in K\)，G 的所有 t 次幂映射在原点附近都有相同的幂级数展开，因为其中任意两个映射都在 0 的某个邻域上重合。将这个幂级数展开记为 \(x^{[t]}\)。

命题 2. (i) \(\psi_{j} \equiv 0 \bmod \deg j\)。

\begin{enumerate}
\def\labelenumi{(\roman{enumi})}
\setcounter{enumi}{1}
\tightlist
\item
  若 \(t\in \mathbf{K}\)，则
\end{enumerate}

\[
x ^ {[ t ]} = \sum_ {i = 1} ^ {\infty} \binom {t} {i} \psi_ {i} (x)\tag{16}
\]

其中右端的形式幂级数由 (i) 知是有意义的。

（记

\[
\binom {t} {i} = \frac {t (t - 1) \dots (t - i + 1)}{i !}
\]

对所有 \(t\in \mathbf{K}.\)

断言 (i) 由 \(\psi_{j}\) 的定义显然成立。

当 t 为整数 \(\geqslant0\) 时，证明 (ii)。由 (14)，

\[
x ^ {[ t ]} = \sum_ {\alpha (1),\ldots,\alpha (t)\in(\mathbf{N}^{n})^{t}} a _ {\alpha (1), \dots , \alpha (t)} x ^ {\alpha (1) + \dots + \alpha (t)}.\tag{17}
\]

对于 \(\alpha = (\alpha(1),\ldots,\alpha(t))\in(\mathbf{N}^{n})^{t}\)，令 \(\sigma(\alpha)\) 表示由 \(j\in\{1,2,\ldots,t\}\) 组成的集合，其中 \(\alpha(j)\neq0\)。若在和式 (17) 中将 \(\sigma(\alpha)\) 相同的项归在一起，则

\[
x ^ {[ t ]} = \sum_ {\sigma \subset (1, t)} h _ {t, \sigma} (x)\tag{18}
\]

其中

\[
h _ {t, \sigma} (x) = \sum_ {\sigma (\alpha) = \sigma} a _ {\alpha (1), \dots , \alpha (t)} x ^ {\alpha (1) + \dots + \alpha (t)}.\tag{19}
\]

令 \(\sigma = \{j_1, j_2, \ldots, j_q\}\)，其中 \(j_1 < j_2 < \ldots < j_q\)。在 (14) 中（将 \(j\) 换成 \(t\)），对 \(x(k)\)，当 \(k \notin \sigma\) 时令其为 0；由于 0 是 G 的单位元，我们得到 \(x(j_1).x(j_2)..x(j_q)\) 在原点附近的幂级数展开：

\[
x (j _ {1}). x (j _ {2}) \dots x (j _ {q}) = \sum_ {\sigma (\alpha) \subset \sigma} a _ {\alpha (1), \dots , \alpha (t)} x (j _ {1}) ^ {\alpha (j _ {1})} x (j _ {2}) ^ {\alpha (j _ {2})} \dots x (j _ {q}) ^ {\alpha (j _ {q})}
\]

从而根据 \(\psi_q\) 的定义：

\[
\psi_ {q} (x) = \sum_ {\sigma (\alpha) = \sigma} a _ {\alpha (1), \dots , \alpha (t)} x ^ {\alpha (j _ {1}) + \dots + \alpha (j _ {q})}.\tag{20}
\]

由 (19) 和 (20) 可见，\(h_{t,\sigma}(x) = \psi_{\mathrm{Card}\sigma}(x)\)。于是 (18) 蕴含

\[
x ^ {[ t ]} = \sum_ {i = 0} ^ {t} \binom {t} {i} \psi_ {i} (x) = \sum_ {i = 0} ^ {\infty} \binom {t} {i} \psi_ {i} (x).
\]

于是记 \(x^{[t]} = \sum_{i=0}^{\infty} \binom{t}{i}\psi_i(x)\) 对所有 \(t\in\mathbf{K}\)。在幂级数 \(x^{[t]}\) 和 \(x^{[t']}\) 中，每个系数都是 \(t\) 的多项式函数。对于 \(x^{[t']}\)，这是显然的。至于 \(x^{[t]}\)，只需证明：对所有 \(u\in\mathrm{U}(G)\)，\(u\) 在映射 \(x\mapsto x^{[t]}\) 下的像是 \(t\) 的多项式函数。现在，对于 \(u\in\mathrm{U}^{m}(G)\)，这个像为 \(t^{m}u\)（§ 4, no. 3, Proposition 7 (iv)）。

由于 \(x^{[t]} = x^{[t']}\) 对整数 \(t\) 满足 \(\geqslant 0\)，于是得到 \(x^{[t]} = x^{[t']}\) 对所有 \(t \in \mathbf{K}\) 都成立。

注记。(1) 当 \(t\) 为整数 \(\geqslant 0\) 时，把命题 2 的条件 (ii) 写成：

\[
\begin{array}{r l} & 0 = \psi_ {0} (x) \\ & x = \psi_ {0} (x) + \psi_ {1} (x) \\ & x ^ {[ 2 ]} = \psi_ {0} (x) + 2 \psi_ {1} (x) + \psi_ {2} (x) \\ & \cdot \cdot \cdot . \end{array}
\]

这些公式足以确定 \(\psi_{i}\)。

\begin{enumerate}
\def\labelenumi{(\arabic{enumi})}
\setcounter{enumi}{1}
\tightlist
\item
  可见 \(\psi_0(x) = 0\)、\(\psi_1(x) = x\)、\(\psi_2(x) = x^{[2]} - 2x\)，
\end{enumerate}

\[
x ^ {[ - 1 ]} = \sum_ {i = 1} ^ {\infty} (- 1) ^ {i} \psi_ {i} (x).
\]

\begin{enumerate}
\def\labelenumi{(\arabic{enumi})}
\setcounter{enumi}{2}
\tightlist
\item
  上述 \(\psi_{2}\) 的表达式和公式 (6) 证明
\end{enumerate}

\[
\psi_ {2} (x) \equiv \mathrm{B} (x, x) \mod \deg 3.\tag{21}
\]

利用命题 2 的 (i) 和 (ii)，可得

\[
x ^ {[ t ]} \equiv t x + \binom {t} {2} \mathrm{B} (x, x) \mod \deg 3.\tag{22}
\]

\begin{enumerate}
\def\labelenumi{(\arabic{enumi})}
\setcounter{enumi}{3}
\tightlist
\item
  令 \(\psi_{p,m}(x)\) 和 \(h_{t,m}(x)\) 表示 \(\psi_p(x)\) 和 \(x^{[t]}\) 的 \(m\) 次齐次分量。于是 \(\psi_{p,m} = 0\) 对 \(m < p\) 成立。另一方面，命题 2 (ii) 给出
\end{enumerate}

\[
h _ {t, m} (x) = \sum_ {p \leqslant m} \frac {t (t - 1) \dots (t - p + 1)}{p !} \psi_ {p, m} (x)\tag{23}
\]

即

\[
h _ {t, m} (x) = \sum_ {1 \leqslant r \leqslant m} t ^ {r} \phi_ {r, m} (x)\tag{24}
\]

其中 \(\phi_{r,m}\) 是从 \(K^{n}\) 到 \(K^{n}\) 的 m 次齐次多项式映射。特别地，由 (23)，

\begin{enumerate}
\def\labelenumi{(\arabic{enumi})}
\setcounter{enumi}{24}
\tightlist
\item
\end{enumerate}

\[
\phi_ {1, m} (x) = \sum_ {p \leqslant m} \frac {(- 1) ^ {p - 1}}{p} \psi_ {p, m} (x)\tag{25}
\]

\[
\phi_ {m, m} (x) = \frac {1}{m !} \psi_ {m, m} (x).\tag{26}
\]

\begin{enumerate}
\def\labelenumi{(\arabic{enumi})}
\setcounter{enumi}{4}
\tightlist
\item
  若 \(\mathbf{K}\) 的特征大于 \(>0\)，则第 1 号和第 2 号的结果仍然成立，条件是第 2 号中的 \(e_{\alpha}\) 定义为 \(\langle e_{\alpha}, \sum_{\beta} \lambda_{\beta} x^{\beta} \rangle = \lambda_{\alpha}\)。在第 3 号中，若按上述方式定义 \(\psi_j\)，则上述论证再次证明 \(x^{[t]} = \sum_{i=1}^{\infty} \binom{t}{i} \psi_i(x)\)，当 \(t \in \mathbf{N}\) 时。
\end{enumerate}

\subsection{指数映射}

令 \(\mathrm{E}(x)\) 为 G 的一个指数映射在 0 附近的幂级数展开。令 \(\mathrm{L}(x)\) 为 G 的一个单射指数映射的逆映射在 0 附近的幂级数展开。由于任一指数映射在 0 处的切映射都是 \(\mathrm{L}(\mathrm{G})\) 的恒等映射，\(\mathrm{E}(x) \equiv x \mod \deg 2\) 且 \(\mathrm{L}(x) \equiv x \mod \deg 2\)。由于 \(\mathrm{E}(\mathrm{L}(x)) = \mathrm{L}(\mathrm{E}(x))\) 对充分接近 0 的 \(x\) 成立，形式幂级数 E 和 L 满足 \(\mathrm{E}(\mathrm{L}(\mathrm{X})) = \mathrm{L}(\mathrm{E}(\mathrm{X})) = \mathrm{X}\)。类似论证表明

\[
\mathrm{E} (t \mathrm{X}) = (\mathrm{E} (\mathrm{X})) ^ {[ t ]}, \quad \mathrm{L} (\mathrm{X} ^ {[ t ]}) = t \mathrm{L} (\mathrm{X})
\]

对 \(t \in K\)。

命题 3.

\begin{enumerate}
\def\labelenumi{(\arabic{enumi})}
\setcounter{enumi}{26}
\tightlist
\item
\end{enumerate}

\[
\mathrm{L} = \sum_ {p = 1} ^ {\infty} \frac {(- 1) ^ {p - 1}}{p} \psi_ {p}\tag{27}
\]

\[
\mathrm{E} = \sum_ {p = 1} ^ {\infty} \frac {1}{p !} \psi_ {p, p}\tag{28}
\]

（回忆 \(\psi_{p,p}\) 是 \(\psi_{p}\) 的 \(p\) 次齐次分量。）

\[
\operatorname{E} (t x) = (\operatorname{E} (x)) ^ {[ t ]}
\]

或者根据 (24)，

\[
\mathrm{E} (t x) = \sum_ {m \geqslant 0} \sum_ {1 \leqslant r \leqslant m} t ^ {r} \phi_ {r, m} (\mathrm{E} (x)).
\]

两边都是关于 \(t\) 和 \(x\) 的形式幂级数。比较 \(t\) 的一次项，得到

\[
x = \sum_ {m \geqslant 0} \phi_ {1, m} (\mathrm{E} (x)).\tag{29}
\]

现在，形式幂级数系统的逆若存在，则唯一存在（Algebra, Chapter IV, § 6, Corollary to Proposition 8）。于是

\[
\begin{array}{r l} \mathrm{L} (x) & = \sum_ {m \geq 0} \phi_ {1, m} (x) \quad \text { by   (29) } \\ & = \sum_ {p, m} \frac {(- 1) ^ {p}}{p} \psi_ {p, m} (x) \quad \text { by   (25) } \\ & = \sum_ {p} \frac {(- 1) ^ {p}}{p} \psi_ {p} (x), \end{array}
\]

从而得到 (i)。类似地，对于 \(t \neq 0\)，

\[
\begin{array}{l} \mathrm{L} (t x) = t \mathrm{L} ((t x) ^ {[ t ^ {- 1} ]}) \\ = t \mathrm{L} \left(\sum_ {m \geqslant 0} \sum_ {1 \leqslant r \leqslant m} t ^ {m - r} \phi_ {r, m} (x)\right) \quad \text { by } (24). \end{array}
\]

比较 t 的一次项，得到

\[
x = \mathrm{L} \Bigl (\sum_ {m \geqslant 0} \phi_ {m, m} (x) \Bigr)
\]

从而

\[
\begin{array}{r l} \mathrm{E} (x) & = \sum_ {m \geqslant 0} \phi_ {m, m} (x) \\ & = \sum_ {m \geqslant 0} \frac {1}{m !} \psi_ {m, m} (x) \quad \text { by   (26) }. \end{array}
\]

命题 4. G 所用的坐标图为规范坐标图的充要条件是：\(\psi_j = 0\) 对 \(j \geqslant 2\) 成立。

由命题 3，这一条件是充分的。设坐标图是规范的，并且 \(\psi_{i} = 0\) 对 \(2 \leqslant i < n\) 成立。则 \(nx = x^{[n]} = \sum_{i=0}^{n} \binom{n}{i} \psi_{i}(x) = nx + \psi_{n}(x)\)，从而 \(\psi_{n} = 0\)。因此，\(\psi_{j} = 0\) 对 \(j \geqslant 2\) 成立，这由关于 \(j\) 的归纳得到。

\section{实李群和复李群}

在本节中，假定 K 等于 R 或 C。

\subsection{从李代数态射到李群态射的过渡}

引理 1. 设 G 是单连通†拓扑群，W 是 e 的对称连通开邻域，H 是一个群，f 是从 W³ 到 H 的映射，使得

\[
f (x y z) = f (x) f (y) f (z)
\]

当 \(x, y, z\) 属于 W 时，存在一个态射 \(f'\)，它从 G 到 H，且满足 \(f' \mid \mathrm{W} = f \mid \mathrm{W}\)。对于 \((g, h) \in \mathrm{G} \times \mathrm{H}\) 以及 W 中 \(e\) 的一个开邻域 U，令 \(\mathrm{A}(g, h, \mathrm{U})\) 为由 \((gu, hf(u)) \in \mathrm{G} \times \mathrm{H}\)（其中 \(u \in \mathrm{U}\)）组成的集合。于是 \((g, h) \in \mathrm{A}(g, h, \mathrm{U})\)，且 \(\mathrm{A}(g, h, \mathrm{U}_1) \cap \mathrm{A}(g, h, \mathrm{U}_2) = \mathrm{A}(g, h, \mathrm{U}_1 \cap \mathrm{U}_2)\)。令 \((s, t) \in \mathrm{A}(g, h, \mathrm{U})\)；有 \(s = gu\) 且 \(t = hf(u)\)，其中 \(u \in \mathrm{U}\)；存在一个开邻域 \(\mathrm{U}'\)，它是 W 中 \(e\) 的邻域，使得 \(u\mathrm{U}' \in \mathrm{U}\)；于是对 \(u' \in \mathrm{U}'\)，

\[
(s u ^ {\prime}, t f (u ^ {\prime})) = (g u u ^ {\prime}, h f (u u ^ {\prime})) \in \mathrm{A} (g, h, \mathrm{U})
\]

从而 \(\mathrm{A}(s,t,\mathrm{U}^{\prime})\subset \mathrm{A}(g,h,\mathrm{U})\)。于是 \(\mathrm{A}(g,h,\mathrm{U})\) 构成 \(\mathbf{G}\times \mathbf{H}\) 上拓扑的一个基。我们把 \(\mathbf{Y}\) 记为带有此拓扑的 \(\mathbf{G}\times \mathbf{H}\)，并把规范投影 \(p\) 记为从 \(\mathbf{Y}\) 到 \(\mathbf{G}\) 的投影，且它是开映射。\(p\) 在 \(\mathrm{A}(g,h,\mathrm{U})\) 上的限制是从 \(\mathrm{A}(g,h,\mathrm{U})\) 到 \(g\mathrm{U}\) 的同胚。因此 \((\mathrm{Y},p)\) 是 \(\mathbf{G}\) 的覆盖空间。令 \(\mathrm{Y}_0\) 为 \(\mathrm{Y}\) 中由 \(\mathrm{A}(e,e,\mathrm{W})\) 生成的子群，令 \(\mathfrak{B}\) 为 \(\mathrm{A}(e,e,\mathrm{U})\) 的集合。显然，\(\mathfrak{B}\) 满足 General Topology, Chapter III, § 1, no. 2 的条件 \((\mathrm{GV}_{\mathrm{I}}^{\prime})\) 和 \((\mathrm{GV}_{\mathrm{II}}^{\prime})\)。集合 \(\mathrm{Y}_0^{\prime}\) 由满足下述条件的 \(y\in \mathrm{Y}_0\) 组成：映射 \(z\mapsto yzy^{-1}\) 和 \(z\mapsto y^{-1}zy\) 从 \(\mathrm{Y}_0\) 到 \(\mathrm{Y}_0\)，并在 \((e,e)\) 处连续；它是 \(\mathrm{Y}_0\) 的一个子群。令 \(w\in W\)。映射 \(w'\mapsto ww'w^{-1}\) 从 \(\mathrm{W}\) 到 \(\mathrm{G}\) 连续，因而映射

\[
\left(w ^ {\prime}, f \left(w ^ {\prime}\right)\right) \mapsto \left(w w ^ {\prime} w ^ {- 1}, f \left(w w ^ {\prime} w ^ {- 1}\right)\right)
\]

从 \(\mathrm{A}(e,e,\mathrm{W})\) 到 \(\mathrm{Y}_0\) 在 \((e,e)\) 处连续。现在

\[
f \left(w w ^ {\prime} w ^ {- 1}\right) = f (w) f \left(w ^ {\prime}\right) f \left(w ^ {- 1}\right)
\]

\[
\left(w w ^ {\prime} w ^ {- 1}, f \left(w w ^ {\prime} w ^ {- 1}\right)\right) = (w, f (w)) \left(w ^ {\prime}, f \left(w ^ {\prime}\right)\right) \left(w, f (w)\right) ^ {- 1}
\]

由于 \(w^{-1} \in W\)，可见 \((w, f(w)) \in Y_0'\)。因此 \(A(e, e, W) \subset Y_0'\)，从而 \(Y_0' = Y_0\)。于是群 \(Y_0\) 带有滤子基 \(\mathfrak{B}\)，并满足 General Topology, Chapter III, § 1, no. 2 的条件 \((GV_{\mathrm{III}}')\)。由于

\[
(g, h). \mathrm{A} (e, e, \mathrm{U}) = \mathrm{A} (g, h, \mathrm{U}),
\]

\(Y_{0}\) 是一个拓扑群，并且由于 \(A(e, e, W)\) 连通而连通。于是 \(p(Y_{0})\) 是 G 的开子群；由于 G 连通，故 \(p(Y_{0}) = G\)。\(p \mid Y_{0}\) 的核是离散的。由于 G 单连通，\(p \mid Y_{0}\) 是从 \(Y_{0}\) 到 G 的同胚。因此，\(Y_{0}\) 是从 G 到 H 的态射 \(f'\) 的图。对于 \(g \in W\)，\((g, f(g)) \in A(e, e, W) \subset Y_{0}\)，从而 \(f(g) = f'(g)\)。

定理 1. 设 G 和 H 是李群，h 是从 L(G) 到 L(H) 的连续态射。假设 G 单连通。则存在唯一一个从 G 到 H 的李群态射 \(\phi\)，使得 \(h = \mathrm{L}(\phi)\)。

\(\phi\) 的存在性由引理 1 和 §4, no. 1, Theorem 1 (i) 得出。\(\phi\) 的唯一性由 §4, no. 1, Theorem 1 (ii) 以及 \(G\) 连通这一事实得出。

推论。设 G 是有限维单连通李群。则存在一个有限维解析线性表示，其核是离散的。

根据 Chapter I, § 7, Theorem 2，存在有限维向量空间 E 以及一个单射态射 \(h\)，它从 \(L(G)\) 到李代数 \(\operatorname{End}(E)\)。由定理 1，存在一个态射 \(\phi\)，它从 G 到 \(\mathbf{GL}(E)\)，使得 \(L(\phi) = h\)。因此 \(\phi\) 是浸入，从而其核是离散的。

注记。(1) 存在没有有限维单射解析线性表示的有限维单连通李群（练习 2）。

\begin{enumerate}
\def\labelenumi{(\arabic{enumi})}
\setcounter{enumi}{1}
\tightlist
\item
  存在有限维连通李群 G，使 G 的每个有限维解析线性表示都有非离散的核（练习 3 和 4）。
\end{enumerate}

\subsection{积分子群}

定义 1. 设 G 是李群。G 的一个积分子群是一个子群 H，带有连通李群结构，且 H 到 G 的规范嵌入是浸入。

\(G\) 的一参数子群是 \(G\) 的一个一维积分子群。

设 H 是 G 的积分子群，i 是 H 到 G 的规范嵌入。则 \(L(i)\) 确定了从 \(L(H)\) 到 \(L(G)\) 中一个具有拓扑补空间的李子代数的同构。以下将 \(L(H)\) 与其在 \(L(i)\) 下的像识别。

例。(1) \(\mathbf{G}\) 的一个连通李子群是 \(\mathbf{G}\) 的一个积分子群。

\begin{enumerate}
\def\labelenumi{(\arabic{enumi})}
\setcounter{enumi}{1}
\item
  假设 \(\mathbf{G}\) 是有限维的。令 \(\mathbf{H}\) 为 \(\mathbf{G}\) 的一个子群；赋予它由 \(\mathbf{G}\) 上的李群结构诱导的结构（§ 4, no. 5, Definition 3）。则其单位元分支 \(\mathrm{H}_0\) 是 \(\mathbf{G}\) 的一个积分子群，且 \(\mathbf{H}\) 在 \(e\) 处的切子代数是 \(\mathrm{L}(\mathrm{H}_0)\)（§ 4, no. 5, Proposition 9 (ii)）。
\item
  设 \(\mathbf{G}\) 是复李群，\(\mathbf{H}\) 是 \(\mathbf{G}\) 的积分子群，\(\mathbf{G}_1\)（分别为 \(\mathbf{H}_1\)）是 \(\mathbf{G}\)（分别为 \(\mathbf{H}\)）的底层实李群。则 \(\mathbf{H}_1\) 是 \(\mathbf{G}_1\) 的积分子群，且 \(\mathbf{L}(\mathbf{H}_1)\) 是 \(\mathbf{L}(\mathbf{H})\) 的底层实李代数。
\end{enumerate}

定理 2. 设 G 是李群。

\begin{enumerate}
\def\labelenumi{(\roman{enumi})}
\item
  映射 \(\mathbf{H} \mapsto \mathbf{L}(\mathbf{H})\) 把 \(\mathbf{G}\) 的积分子群集合双射到 \(\mathbf{L}(\mathbf{G})\) 中具有拓扑补空间的李子代数集合。
\item
  设 \(\mathbf{H}\) 是 \(\mathbf{G}\) 的积分子群。\(\mathbf{G}\) 的每个连通李子群芽，其李代数为 \(\mathbf{L}(\mathbf{H})\)，都是 \(\mathbf{H}\) 的开子流形，并生成 \(\mathbf{H}\)。
\end{enumerate}

\begin{enumerate}
\def\labelenumi{(\alph{enumi})}
\item
  令 \(\mathfrak{b}\) 为 \(\mathrm{L}(\mathrm{G})\) 中具有拓扑补空间的李子代数。令 \(\mathrm{H}_{1}\) 为 \(\mathrm{G}\) 的李子群芽，使得 \(\mathrm{L}(\mathrm{H}_{1}) = \mathfrak{b}\)（§ 4, Theorem 3）。可以选择 \(\mathrm{H}_{1}\) 使其连通。令 \(\mathrm{H}\) 为 \(\mathrm{G}\) 中由 \(\mathrm{H}_{1}\) 生成的子群。存在一个 \(\mathrm{H}\) 上的李群结构（§ 1, Corollary to Proposition 22），使得 \(\mathrm{H}_{1}\) 是 \(\mathrm{H}\) 的开子流形，且 \(\mathrm{H}\) 到 \(\mathrm{G}\) 的规范嵌入是浸入。由于 \(\mathrm{H}_{1}\) 连通，\(\mathrm{H}\) 连通，因而是 \(\mathrm{G}\) 的积分子群。于是 \(\mathrm{L}(\mathrm{H}) = \mathrm{L}(\mathrm{H}_{1}) = \mathfrak{b}\)。这证明了 (i) 中所考虑的映射是满射。
\item
  设 H 是 G 的积分子群，\(N_{1}\) 是李代数为 L(H) 的 G 的连通李子群芽。由于 H 到 G 的规范嵌入是浸入，存在 H 的一个开子群芽 \(H_{1}\)，它同时是 G 的子流形，因而是李代数为 L(H) 的 G 的李子群芽。另一方面，令 N 为由 \(N_{1}\) 生成的 G 的子群；由证明的 (a) 部分，N 带有 G 的积分子群结构，使 \(N_{1}\) 是 N 的开子流形。由 §4, Theorem 3，\(H_{1} \cap N_{1}\) 在 \(H_{1}\) 和 \(N_{1}\) 中都是开集。因此，由 \(H_{1} \cap N_{1}\) 生成的 G 的子群一方面等于 H，另一方面等于 N。所以李群 H 和 N 相等。这证明了 (ii)，也证明了 (i) 中所考虑的映射是单射。
\end{enumerate}

注记 1. 设 H 是 G 的积分子群。令 Y 为与 L(H) 相关的 G 的左叶层。若 \(g \in G\)，则通过 \(\gamma(g)\) 赋予 gH 由 H 的流形结构导出的流形结构。由 Differentiable and Analytic Manifolds, R, 9.3.2，gH 到 Y 的规范嵌入是态射。这个态射是 étale 的。因此，Y 的极大连通叶是模 H 的左陪集。

命题 1. 设 G 和 M 是李群，H 是 G 的积分子群，\(\phi\) 是从 M 到 G 的态射，且 \(\mathrm{L}(\phi)(\mathrm{L}(\mathrm{M})) \subset \mathrm{L}(\mathrm{H})\)。假设 M 连通。则 \(\phi(\mathrm{M}) \subset \mathrm{H}\)，并且把 \(\phi\) 看作从 M 到 H 的映射时，它是李群态射。

在注记 1 的记号下，\(\phi\) 是从 M 到 Y 的态射（Differentiable and Analytic Manifolds, R, 9.3.2），由于 M 连通，故 \(\phi(M) \subset H\)。

推论 1. 设 G 和 H 是李群，\(\phi\) 是从 G 到 H 的李群态射，N 是 \(\phi\) 的核，且 \(h = L(\phi)\)。假设 G 连通且 H 有限维。

\begin{enumerate}
\def\labelenumi{(\roman{enumi})}
\item
  N 是 G 的李子群，且 L(N) = Ker h。
\item
  令 \(\mathbf{H}'\) 为 \(\mathbf{H}\) 的积分子群，其李代数为 \(\operatorname{Im} h\)。则 \(\phi(\mathbf{G}) = \mathbf{H}'\)。
\item
  将 \(\phi\) 通过商群得到的从 G/N 到 H' 的映射是李群同构。
\end{enumerate}

\((i)\) 已经证明（§ 3, no. 8, Proposition 28）。

令 \(\psi\) 为将 \(\phi\) 通过商群得到的从 G/N 到 H 的李群态射；它是浸入（§3, no. 8, Proposition 28）。由命题 1，\(\psi\) 是从 G/N 到 H' 的李群态射。这个态射是 étale 的，由于 H' 连通，故 \(\psi\) (G/N) = H'；这证明了 (ii)。于是 \(\psi\colon G/N \to H'\) 是双射，也是李群同构，从而证明 (iii)。

推论 2. 设 \(\mathbf{G}\) 是李群，\(\mathrm{H}_1\) 和 \(\mathrm{H}_2\) 是 \(\mathbf{G}\) 的积分子群。若 \(\mathrm{L}(\mathrm{H}_2) \subset \mathrm{L}(\mathrm{H}_1)\)，则 \(\mathrm{H}_2\) 是 \(\mathrm{H}_1\) 的积分子群。

令 \(i_1: \mathrm{H}_1 \to \mathrm{G}\)、\(i_2: \mathrm{H}_2 \to \mathrm{G}\) 为规范嵌入。于是

\[
\mathrm{L} (i _ {2}) (\mathrm{L} (\mathrm{H} _ {2})) = \mathrm{L} (\mathrm{H} _ {2}) \subset \mathrm{L} (\mathrm{H} _ {1}).
\]

由命题 1，\(i_{2}\) 是从 \(H_{2}\) 到 \(H_{1}\) 的解析映射，甚至是从 \(H_{2}\) 到 \(H_{1}\) 的浸入；这是因为 \(\mathrm{L}(i_{2})\) 是从 \(\mathrm{L}(\mathrm{H}_{2})\) 到 \(\mathrm{L}(\mathrm{H}_{1})\) 的某个具有拓扑补空间的子代数的同构。

推论 3. 设 \(\mathbf{G}\) 是有限维李群，\((\mathrm{H}_i)_{i\in I}\) 是 \(\mathbf{G}\) 的一族李子群。则 \(\mathrm{H} = \bigcap_{i\in I}\mathrm{H}_i\) 是 \(\mathbf{G}\) 的李子群，并且

\[
\mathrm{L} (\mathrm{H}) = \bigcap_ {i \in \mathrm{I}} \mathrm{L} (\mathrm{H} _ {i}).
\]

存在 I 的一个有限子集 J，使得 \(\bigcap_{i\in J}L(H_i)\) 等于所有 \(\mathrm{L}(\mathrm{H}_i)\) 的交 M。我们知道 \(\mathrm{H}^* = \bigcap_{i\in J}\mathrm{H}_i\) 是满足 \(\mathrm{L}(\mathrm{H}^*) = \mathrm{M}\) 的李子群（§ 3, no. 8, Corollary 2 to Proposition 29）。令 \(\mathrm{H}_0\) 为 \(\mathrm{H}^*\) 的单位元分支。它是 G 的李子群，且 \(\mathrm{L}(\mathrm{H}_0) = \mathrm{M}\)。由推论 2，\(\mathrm{H}_0\subset \mathrm{H}_i\) 对所有 \(i\) 成立，因而 \(\mathrm{H}_0\subset \mathrm{H}\subset \mathrm{H}^*\)，从而得到推论。

推论 4. 设 G 是有限维连通李群。以下条件等价：

\begin{enumerate}
\def\labelenumi{(\roman{enumi})}
\item
  G 是幺模的（Integration, Chapter VII, § 1, no. 3, Definition 3）；
\item
  对所有 g ∈ G，都有 det Ad g = 1；
\item
  \(\operatorname{Tr} \operatorname{ad} a = 0\) 对所有 \(a \in \mathrm{L}(\mathrm{G})\) 都成立。
\end{enumerate}

映射 \(g \mapsto \operatorname{det} \operatorname{Ad} g\) 是一个态射 \(\phi\)，它从 \(G\) 到 \(K^*\)。由 § 3, Propositions 35 (no. 10) and 44 (no. 12)，有 \(L(\phi)a = Tr\operatorname{ad} a\) 对所有 \(a \in L(G)\) 成立。显然 \(\operatorname{Im} L(\phi) = \{0\}\) 或 \(K\)。在第一种（分别为第二种）情形，由推论 1，\(\operatorname{Im} \phi = \{1\}\)（分别为 \(\operatorname{Im} \phi = K^*\)），因而根据 § 3, no. 16, Corollary to Proposition 55，\(G\) 是幺模的（分别为非幺模的）。

命题 2. 设 \(\mathbf{G}\) 是有限维李群，\(\mathbf{H}\) 是 \(\mathbf{G}\) 的积分子群。以下条件等价：

\begin{enumerate}
\def\labelenumi{(\roman{enumi})}
\item
  H 是闭的；
\item
  H 上的拓扑由 G 上的拓扑诱导；
\item
  H 是 G 的李子群。
\end{enumerate}

\((i) \Rightarrow (iii)\)：这由 § 1, Propositions 2 (iv) (no. 1) and 14 (iii) (no. 7) 得出。

\((iii) \Rightarrow (ii)\)：显然。

\((ii) \Rightarrow (i)\)：若 H 上的拓扑由 G 上的拓扑诱导，则 H 是闭的，因为 H 是完备的（§ 1, no. 1, Proposition 1）。

命题 3. 设 G 是李群，H 是 G 的积分子群，M 是非空连通解析流形，f 是从 M 到 G 的映射，且 \(r \in N_{K}\)。考虑以下条件：

\begin{enumerate}
\def\labelenumi{(\roman{enumi})}
\item
  \(f\) 属于 \(\mathbf{C}^r\) 类，且 \(f(\mathbf{M})\subset \mathbf{H}\)；
\item
  \(f(\mathbf{M})\subset \mathbf{H}\)，并且把 \(f\) 看作从 \(\mathbf{M}\) 到 \(\mathbf{H}\) 的映射时，属于 \(\mathbf{C}^r\) 类
\item
  \(f\) 属于 \(\mathbf{C}^r\) 类，\(f(\mathbf{M})\) 与 \(\mathbf{H}\) 相交，并且 \(\mathrm{T}_m(\mathbf{M})\) 的像包含于 \(f(m)\) . \(\mathbf{L}(\mathbf{H})\)，对所有 \(m \in \mathbf{M}\)。
\end{enumerate}

则 (ii) \(\Leftrightarrow\) (iii) \(\Rightarrow\) (i)。若 H 上的拓扑有可数基，则三个条件等价。

\((ii) \Rightarrow (i)\) 以及 \((ii) \Rightarrow (iii)\)：显然。

\((iii) \Rightarrow (ii)\)：设条件 (iii) 成立。由 Differentiable and Analytic Manifolds, R, 9.2.8，\(f\) 是一个 \(\mathbf{C}^r\) 类态射，从 M 到与 \(\mathbf{L}(\mathbf{H})\) 相关的左叶层。由于 M 连通，\(f(\mathbf{M}) \subset \mathbf{H}\)。

若 H 上的拓扑有可数基，则条件 (i) 蕴含 f 是从 M 到 H 的 \(C^{r}\) 类映射（Differentiable and Analytic Manifolds, R, 9.2.8）；因此 (i) \(\Rightarrow\) (ii)。

推论 1. 设 G 是有限维李群，H 是 G 的积分子群。则 H 在 e 处的切李子代数（§ 4, no. 5, Definitions 2 and 3）是 L(H)，且 H 上的李群结构是由 G 上的结构诱导的结构。

由于 H 连通且有限维，其拓扑有可数基。

推论 2. 设 G 是李群，\(H_{1}\) 和 \(H_{2}\) 是 G 的积分子群。假设 \(H_{1}\) 上的拓扑有可数基。则

\[
\mathrm{H} _ {2} \subset \mathrm{H} _ {1} \Leftrightarrow \mathrm{L} (\mathrm{H} _ {2}) \subset \mathrm{L} (\mathrm{H} _ {1})
\]

并且，若这些条件成立，\(\mathbf{H}_2\) 是 \(\mathbf{H}_1\) 的积分子群。

最后一个断言以及蕴含 \(\mathrm{L}(\mathrm{H}_{2})\subset\mathrm{L}(\mathrm{H}_{1})\Rightarrow\mathrm{H}_{2}\subset\mathrm{H}_{1}\) 由 Proposition 1 的 Corollary 2 得出。逆向蕴含由命题 3 得出。

推论 3. 设 G 是李群，\(H_{1}\) 和 \(H_{2}\) 是 G 的积分子群，且它们的拓扑有可数基。若 \(H_{1}\) 和 \(H_{2}\) 具有相同的底层集合，则 \(H_{1}\) 和 \(H_{2}\) 上的李群结构相同。

这由推论 2 得出。

注记 2. 设 G 是有限维李群，H 是 G 的子群。若 H 上存在李群结构 S，使得带有 S 的 H 是 G 的积分子群，则我们约定称 H 是 G 的积分子群，这是一种语言上的滥用。由 Proposition 3 的 Corollary 3，若 S 存在，则 S 唯一。

注记 3. 设 V 是 \(\mathbf{C}^r\) 类流形。令 M 为 V 的子集，\(x\) 和 \(y\) 为 M 的元素。考虑以下性质：

\(P_{M,x,y}\) ：存在 I、\(x_{0}, x_{1}, \ldots, x_{n}, f_{1}, \ldots, f_{n}\)，使得：(a) I 是 K 的连通开子集；(b) \(x_{0}, \ldots, x_{n}\) 属于 M，\(x_{0} = x, x_{n} = y\)；(c) 对 \(1 \leqslant i \leqslant n\)，\(f_{i}\) 是从 I 到 V 的 \(C^{r}\) 类映射，取值为 \(x_{i-1}\) 和 \(x_{i}\)，且 \(f_{i}(I) \subset M\)。

称 M 为 V 的 \(C^{r}\)-连通子集，当且仅当对 M 中所有元素 x, y，性质 \(P_{M,x,y}\) 都成立。

命题 4. 设 G 是有限维李群，H 是 G 的子群。令 \(r \in \mathbf{N}_{\mathbb{K}}\)。以下条件等价：

\begin{enumerate}
\def\labelenumi{(\roman{enumi})}
\item
  H 是 G 的积分子群；
\item
  赋予 H 由 G 上的李群结构诱导的李群结构后，H 是连通的；
\item
  H 是 \(\mathbf{C}^r\)-连通的。
\end{enumerate}

\((ii) \Rightarrow (i)\)：显然。

\((i) \Rightarrow (iii)\)：设 H 带有李群结构，使 H 成为 G 的积分子群。使用注记 3 的记号，\(y \in H\) 中满足性质 \(\mathrm{P}_{\mathrm{H},e,y}\) 的元素所成的集合是 H 的一个开子群。由于 H 连通，这个子群等于 H，因而条件 (iii) 成立。

\((iii) \Rightarrow (ii)\)：设条件 (iii) 成立，并赋予 H 由 G 上的李群结构诱导的结构。令 \(\mathfrak{h}\) 为 H 在 \(e\) 处的切子代数。H 的单位元分支 \(\mathrm{H}_0\) 是 G 的积分子群，且 \(\mathrm{L}(\mathrm{H}_0) = \mathfrak{h}\)。我们证明 \(\mathrm{H} = \mathrm{H}_0\)。只需证明如下事实：令 I 为 K 的连通开子集，\(f\) 为从 I 到 G 的 \(\mathbf{C}^r\) 类映射，使 \(f(\mathrm{I}) \subset \mathrm{H}\)，令 \(\lambda\) 和 \(\mu\) 为 I 中两点；若 \(f(\lambda) \in \mathrm{H}_0\)，则 \(f(\mu) \in \mathrm{H}_0\)。但是，对所有 \(\nu \in \mathrm{I}\)，都有 \((\mathrm{T}_{\nu}f)(\mathrm{K}) \subset f(\nu)\mathfrak{h}\)，这是由 \(\mathfrak{h}\) 的定义得出的，所以我们的断言由命题 3 得出。

注记 4. 若 \(\mathbf{K} = \mathbf{R}\)，则 \(\mathbf{G}\) 的积分子群还可以刻画为这样的子群：赋予它们由 \(\mathbf{G}\) 的拓扑诱导的拓扑后，它们是弧连通的（§ 8, Exercise 4）。然而，子群可能是连通的，却不是积分子群（Commutative Algebra, Chapter VI, § 9, Exercise 2）。

推论。设 G 是有限维李群，\(H_{1}\) 和 \(H_{2}\) 是 G 的两个积分子群。由 \(H_{1}\) 和 \(H_{2}\) 生成的 G 的子群以及 G 的子群 ( \(H_{1}, H_{2}\) ) 都是 G 的积分子群。

G 的子群 (G, G) 不总是闭的（§ 9, Exercise 6）。

回忆（§ 3, no. 11, Corollary 5 to Proposition 41）：若 a 是有限维代数，则 \(\operatorname{Aut}(a)\) 是 \(\mathbf{GL}(a)\) 的李子群，且 \(\mathrm{L}(\mathrm{Aut}(a))\) 是 \(a\) 的导子李代数。

定义 2. 设 \(\mathfrak{a}\) 是有限维李代数。令 \(\operatorname{Ad}(\mathfrak{a})\) 或 \(\operatorname{Int}(\mathfrak{a})\) 表示 \(\operatorname{Aut}(\mathfrak{a})\) 中李代数为 \(\operatorname{ad}(\mathfrak{a})\) 的积分子群。这个群的元素称为 \(\mathfrak{a}\) 的内自同构。

由结构传递，\(\operatorname{ad}(a)\) 在 \(\operatorname{Aut}(a)\) 下不变，因而 \(\operatorname{Int}(a)\) 是 \(\operatorname{Aut}(a)\) 的正规子群。由 § 4, no. 4, Corollary 1 to Proposition 8 以及 \(\operatorname{Int}(a)\) 连通这一事实，\(\operatorname{Int}(a)\) 的元素是形如 \(\exp ad x\)（其中 \(x \in a\)）的自同构的有限乘积。一般而言，\(\operatorname{Int}(a)\) 不是 \(\operatorname{Aut}(a)\) 的李子群（练习 14）。

\subsection{从李代数到李群的过渡}

定理 3. (i) 若 L 是有限维李代数，则存在一个单连通李群 G，使 L(G) 同构于 L。

\begin{enumerate}
\def\labelenumi{(\roman{enumi})}
\setcounter{enumi}{1}
\tightlist
\item
  设 \(G_{1}\) 和 \(G_{2}\) 是两个连通李群，且 \(G_{1}\) 单连通。令 f 为从 \(\mathrm{L}(\mathrm{G}_{1})\) 到 \(\mathrm{L}(\mathrm{G}_{2})\) 的同构，\(\phi\) 为从 \(G_{1}\) 到 \(G_{2}\) 的态射，满足 \(\mathrm{L}(\phi)=f\)，N 为 \(\phi\) 的核。则 N 是 \(G_{1}\) 的中心中的离散子群，而从 \(G_{1}/N\) 到 \(G_{2}\) 的态射由 \(\phi\) 通过商得到，且是李群同构。若 \(G_{2}\) 单连通，则 \(\phi\) 是同构。
\end{enumerate}

令 \(\mathbf{L}\) 为有限维李代数。存在有限维向量空间 \(\mathbf{E}\)，使得 \(\mathbf{L}\) 可以识别为 \(\operatorname{End}(\mathbf{E})\) 的李子代数（Chapter I, §7, Theorem 2）。令 \(\mathbf{H}\) 为 \(\mathbf{GL}(\mathbf{E})\) 中李代数为

L 的积分子群。令 \(\hat{\mathbf{H}}\) 为它的万有覆盖（§ 1, no. 9, Remark）。则 \(\mathbf{L}(\hat{\mathbf{H}})\) 同构于 L，从而得到 (i)。

  令 \(G_{1}, G_{2}, \tilde{f}, \phi, N\) 如 (ii) 中所述。则 \(\phi\) 是 étale 的，因而 \(\phi(G_{1})\) 是 \(G_{2}\) 的开子群；于是 \(\phi(G_{1}) = G_{2}\)。另一方面，N 是离散的，因而包含于 \(G_{1}\) 的中心中（Integration, Chapter VII, § 3, Lemma 4）。显然，从 \(G_{1}/N\) 到 \(G_{2}\) 的态射由 \(\phi\) 通过商得到，且是李群同构。若 \(G_{2}\) 单连通，则每个从 \(G_{1}\) 到 \(G_{2}\) 的 étale 映射都是单射，因而 \(N = \{e\}\)。

命题 5. 设 G 是连通实李群。假设 L(G) 带有一个复可赋范李代数结构 L'，与其可赋范实李代数结构相容。G 上存在唯一一个复李群结构，该结构与实李群结构相容，且李代数为 L'。

由 § 4, no. 2, Corollary 2 to Theorem 2，只需证明 \(L'\) 的结构在 Ad G 下不变。令 \(\phi\) 为 G 的一个指数映射。由 § 4, no. 4, Corollary 3 (i) to Proposition 8，存在 \(L(G)\) 中 0 的一个邻域 V，使 \(L'\) 的结构在 \(Ad\phi(V)\) 下不变。但是 \(Ad\phi(V)\) 生成 G，因为 G 是连通的。

若不假设 G 连通，命题 5 的结论不一定成立（练习 7）。

命题 6. 设 G 是连通复李群。若 G 紧，则 G 交换。

全纯映射 \(g \mapsto \operatorname{Ad} g\) 从 \(\mathbf{G}\) 到 \(\mathcal{L}(\mathbf{L}(\mathbf{G}))\) 是常值的（Differentiable and Analytic Manifolds, R, 3.3.7），因而 ad \(a = 0\) 对所有 \(a \in \mathbf{L}(\mathbf{G})\) 成立（§ 3, no. 12, Proposition 44）。所以 \(\mathbf{G}\) 是交换的（§ 4, Corollary 3 to Theorem 1）。

\subsection{指数映射}

定理 4. 设 \(\mathbf{G}\) 是李群。存在唯一一个 \(\mathbf{G}\) 的指数映射，定义在 \(\mathrm{L}(\mathbf{G})\) 上。

存在一个凸开邻域 \(\mathbf{U}\)，它位于 \(\mathrm{L}(\mathbf{G})\) 中并包含 0，以及一个 \(\mathbf{G}\) 的指数映射，定义在 \(\mathbf{U}\) 上。通过把 \(\mathbf{U}\) 取得充分小，可以假设

\[
\phi ((\lambda + \lambda^ {\prime}) a) = \phi (\lambda a) \phi (\lambda^ {\prime} a)\tag{1}
\]

其中 \(a \in \mathrm{L}(\mathbf{G})\)，\(\lambda, \lambda'\) 属于 \(\mathbf{K}\)，且 \(\lambda a, \lambda'a\)、\((\lambda + \lambda')a\) 属于 \(\mathbf{U}\)。

令 \(a \in \mathrm{L}(\mathrm{G})\)。存在一个整数 \(n > 0\)，使 \(\frac{1}{n} a \in \mathrm{U}\)。若 \(m\) 是满足 \(\frac{1}{m} a \in \mathrm{U}\) 的另一个整数，则 \(\frac{1}{nm} a \in \mathrm{U}\)，且关系 (1) 蕴含

\[
\phi \left(\frac {1}{n} a\right) = \left(\phi \left(\frac {1}{n m} a\right)\right) ^ {m}, \quad \phi \left(\frac {1}{m} a\right) = \left(\phi \left(\frac {1}{n m} a\right)\right) ^ {n};
\]

从而 \(\left(\phi\left(\frac{1}{n}a\right)\right)^n = \left(\phi\left(\frac{1}{m}a\right)\right)^m\)。存在一个延拓 \(\psi: \mathrm{L}(G) \to G\)，它是 \(\phi\) 的延拓，并满足 \(\psi(a) = \left(\phi\left(\frac{1}{n}a\right)\right)^n\) 对 \(a \in \mathrm{L}(G)\) 以及 \(n\) 满足 \(>0\) 且 \(\frac{1}{n}a \in \mathrm{U}\) 时成立。显然 \(\psi\) 是解析的，也是 \(G\) 的指数映射。若 \(\psi': \mathrm{L}(G) \to G\) 是 \(G\) 的一个指数映射，则 \(\psi\) 和 \(\psi'\) 在 0 的某个邻域上重合；由于 \(\mathrm{L}(G)\) 连通，故二者相等。

今后提到 G 的指数映射时，均指定理 4 中所考虑的映射。若不会引起混淆，将其记为 \(\exp_{G}\) 或 exp。

例。设 A 是完备赋范含幺结合代数。则 \(\exp_{\mathbf{A}}\) 是 Chapter II, § 7, no. 3 中定义的指数映射。

命题 7. 设 G 是李群，a 是 L(G) 的元素。从 K 到 G 的映射 \(\lambda \mapsto \exp(\lambda a)\) 是唯一的李群 K 到 G 的态射 \(\phi\)，满足 \((\mathrm{T}_0\phi)1 = a\)。

映射 \((\lambda, \lambda') \mapsto \exp(\lambda a) \exp(\lambda' a)\) 和 \((\lambda, \lambda') \mapsto \exp(\lambda + \lambda') a\) 从 \(K \times K\) 到 \(G\)，都是解析的，并在 \((0, 0)\) 的某个邻域上重合。由于 \(K \times K\) 连通，这两个映射相等。因此 \(\phi: \lambda \mapsto \exp(\lambda a)\) 是从 \(K\) 到 \(G\) 的李群态射。从 \(0\) 处到映射 \(\lambda \mapsto \lambda a\) 的切映射就是映射 \(\lambda \mapsto \lambda a\)；且 \(T_0(\exp) = \mathrm{Id}_{\mathrm{L}(\mathrm{G})}\)；所以 \((T_0\phi)1 = a\)。命题的唯一性断言由定理 1 得出。

命题 8. 设 G 是李群。对于 L(G) 中所有 x, y 以及整数 n，

\begin{enumerate}
\def\labelenumi{(\arabic{enumi})}
\setcounter{enumi}{1}
\tightlist
\item
\end{enumerate}

\[
\exp (x + y) = \lim _ {n \rightarrow + \infty} \left(\left(\exp \frac {1}{n} x\right)\left(\exp \frac {1}{n} y\right)\right) ^ {n}\tag{2}
\]

\[
\exp [ x, y ] = \lim _ {n \rightarrow + \infty} \left(\left(\exp \frac {1}{n} x\right)\left(\exp \frac {1}{n} y\right)\left(\exp \frac {1}{n} x\right) ^ {- 1} \left(\exp \frac {1}{n} y\right) ^ {- 1}\right) ^ {n ^ {2}}\tag{3}.
\]

由命题 7，取 \(\lambda = \frac{1}{n}\)，这由 § 4, no. 3 的 Proposition 4 得出。

命题 9. 设 G 是复李群，G' 是 G 的底层实李群。则 \(\exp_{G} = \exp_{G'}\)。

这由 §4, no. 3 的 Proposition 5 以及 \(\exp_{G}\) 和 \(\exp_{G'}\) 的解析性得出。

命题 10. 设 G 和 H 是李群，\(\phi\) 是从 G 到 H 的态射。

\begin{enumerate}
\def\labelenumi{(\roman{enumi})}
\item
  \(\phi \circ \exp_{\mathrm{G}} = \exp_{\mathrm{H}} \circ \mathrm{L}(\phi)\) .
\item
  若 G 是 H 的积分子群，则 \(\exp_{G} = \exp_{H} | L(G)\)。
\end{enumerate}

等式 (i) 两边都是从 L(G) 到 H 的解析映射，并在 0 的某个邻域上重合（§ 4, Proposition 8, no. 4），因而相等。断言 (ii) 是 (i) 的特殊情形。

推论 1. 设 G 是李群，G' 是 G 的李子群，且 \(a \in \mathrm{L}(\mathrm{G})\)。以下条件等价：

\begin{enumerate}
\def\labelenumi{(\roman{enumi})}
\item
  \(a \in \mathbf{L}(\mathbf{G}')\) ;
\item
  当 \(\exp (\lambda a)\in \mathbf{G}'\) 对所有 \(\lambda \in \mathbf{K}\) 且 \(|\lambda |\) 充分小时成立；
\item
  \(\exp(\lambda a) \in \mathbf{G}'\) 对所有 \(\lambda \in K\) 成立。
\end{enumerate}

论证与 § 4, no. 4, Corollary 1 to Proposition 8 中的相同。

推论 2. 设 \(\mathbf{G}\) 是李群，\(\mathbf{H}\) 是 \(\mathbf{G}\) 的积分子群，且 \(a \in \mathbf{L}(\mathbf{G})\)。考虑以下条件：

\begin{enumerate}
\def\labelenumi{(\roman{enumi})}
\item
  \(a \in \mathrm{L}(\mathrm{H})\) ;
\item
  \(\exp_{\mathbf{G}}(\lambda a)\in \mathbf{H}\) 对所有 \(\lambda \in \mathbf{K}\) 成立。
\end{enumerate}

则 (i) \(\Rightarrow\) (ii)。若 H 的拓扑有可数基，则 (i) \(\Leftrightarrow\) (ii)。

令 \(i\) 为从 \(\mathbf{H}\) 到 \(\mathbf{G}\) 的规范嵌入。若 \(a \in \mathbf{L}(\mathbf{H})\)，则

\[
\exp_ {\mathrm{G}} (\lambda a) = \left(\exp_ {\mathrm{G}} \circ \mathrm{L} (i)\right) (\lambda a) = (i \circ \exp_ {\mathrm{H}}) (\lambda a) \in \mathrm{H}.
\]

因此 (i) \(\Rightarrow\) (ii)。H 有可数基时的逆向蕴含由命题 3 得出。

推论 3. 设 G 是李群，\(\rho\) 是 G 的解析线性表示，\(x \in L(G)\)，且 \(g \in G\)。

\begin{enumerate}
\def\labelenumi{(\roman{enumi})}
\item
  \(\rho(\exp x) = \exp L(\rho)x;\)
\item
  \(\mathrm{Ad}(\exp x) = \exp \mathrm{ad} x;\)
\item
  \(g(\exp x)g^{-1} = \exp(\operatorname{Ad} g.x).\)
\end{enumerate}

论证与 § 4, no. 4, Corollaries 2 and 3 to Proposition 8 中的相同。

推论 4. 设 G 是有限维连通李群。

\begin{enumerate}
\def\labelenumi{(\roman{enumi})}
\item
  \(\text{Int}(L(G)) = \text{Ad}(G).\)
\item
  令 \(\mathbf{Z}\) 为 \(\mathbf{G}\) 的中心。则 \(\mathbf{Z}\) 是 \(\mathbf{G}\) 的李子群，其李代数是 \(\mathbf{L}(\mathbf{G})\) 的中心。从 \(\mathbf{G} / \mathbf{Z}\) 到 \(\operatorname{Int} \mathbf{L}(\mathbf{G})\) 的映射由 \(g \mapsto \operatorname{Ad} g\) 通过商得到，且是李群同构。
\end{enumerate}

断言 (i) 由推论 3 (ii) 以及定义 2 后的注记得出。令 \(g \in G\)。要使 \(\operatorname{Ad} g = \operatorname{Id}_{\operatorname{L(G)}}\)，充要条件是 \(\operatorname{Int} g\) 与 \(\operatorname{Id}_G\) 在 \(e\) 的某个邻域上重合（§ 4, no. 1, Theorem 1 (ii)），因而在整个 \(G\) 上重合；换言之，充要条件是 \(g \in Z\)。于是 (ii) 由 Proposition 1 的 Corollary 1 得出。

定义 3. 设 G 是有限维连通李群。李群 \(\operatorname{Int}(\mathbf{L}(\mathbf{G})) = \operatorname{Ad}(\mathbf{G})\) 称为 G 的伴随群。

命题 11. 设 \(\mathbf{G}\) 是连通交换李群。

\begin{enumerate}
\def\labelenumi{(\roman{enumi})}
\item
  exp 是从加法李群 L(G) 到 G 的 étale 态射。
\item
  若 \(\mathbf{K} = \mathbf{R}\) 且 \(\mathbf{G}\) 有限维，则 \(\mathbf{G}\) 同构于形如 \(\mathbf{R}^p \times \mathbf{T}^q\) 的李群（\(p, q\) 为整数 \(\geqslant 0\)）。
\end{enumerate}

由 Hausdorff 公式，有 \((\exp x)(\exp y) = \exp (x + y)\)，当 \(x,y\) 充分接近 0 时成立；由解析延拓可知，该等式对所有 \(x,y\) 属于 \(\mathrm{L}(G)\) 时成立。因此 exp 是群同态，并且是 étale 的，因为

\[
\mathrm{T} _ {e} (\exp) = \mathrm{Id} _ {\mathrm{L(G)}}.
\]

从而得到 (i)。断言 (ii) 由 (i) 和 General Topology, Chapter VII, § 1, Proposition 9 得出。

命题 12. 设 G 是李群，L = L(G)。对于所有 \(x \in L\)，将 \(\mathrm{T}_{x}(\mathrm{L})\) 与 L 识别，从而 exp 在 x 处的右微分 \(\varpi(x)\) 是从 L 到 L 的线性映射。对于所有 \(x \in L\)，

\[
\varpi (x) = \sum_ {n \geqslant 0} \frac {1}{(n + 1) !} (\mathrm{ad} x) ^ {n}.
\]

两边都是 \(x\) 的解析函数，并且当 \(x\) 充分接近 0 时相等（§ 4, no. 3, Proposition 6）。

注记。\(\varpi(x) . (\mathrm{ad} x) = \exp \mathrm{ad} x - 1\)。我们滥用记号写成

\[
\varpi (x) = \frac {\exp \operatorname{ad} x - 1}{\operatorname{ad} x}.
\]

推论。设 \(\mathbf{G}\) 是复李群，\(x \in \mathrm{L}(\mathbf{G})\)。在 \(x\) 处，\(\exp_{\mathbf{G}}\) 的切映射的核为 \(\bigoplus_{n \in \mathbf{Z}\setminus\{0\}} \operatorname{Ker}(\operatorname{ad} x - 2i\pi n)\)。

整函数 \(z \mapsto \sum_{n \geqslant 0} \frac{1}{(n + 1)!} z^n\) 等于 \(\frac{e^z - 1}{z}\)（当 \(z \neq 0\) 时），并且具有

零点 \(2\pi i\mathbf{Z}\setminus\{0\}\) 中的各点，且这些都是单零点。于是推论由命题 12 和以下引理得出：

引理 2. 设 E 是复 Banach 空间，u 是 \(\mathcal{L}(\mathrm{E})\) 的元素，S 是 \(u\) 在 \(\mathcal{L}(\mathrm{E})\) 中的谱（Spectral Theories, Chapter I, §1, no. 2），\(f\) 是 S 的一个开邻域 \(\Omega\) 上的全纯复函数。假设 \(f\) 在 \(\Omega\) 中只有有限个互异零点 \(z_{1},\ldots ,z_{n}\)，其重数分别为 \(h_1,\dots,h_n\)。则 \(\operatorname {Ker}f(u)\) 是各 \(\operatorname {Ker}(u - z_i)^{h_i}\)（\(1\leqslant i\leqslant n\)）的直和。

（关于 \(f(u)\) 的定义，见 Spectral Theories, Chapter I, § 4, no. 8。）

存在处处非零的全纯函数 \(g\)，它定义在 \(\Omega\) 上，使得 \(f(z) = (z - z_1)^{h_1} \ldots (z - z_n)^{h_n} g(z)\)。于是 \(g(u)g^{-1}(u) = g^{-1}(u)g(u) = 1\)，从而 \(\operatorname{Ker}f(u) = \operatorname{Ker}\prod_{i=1}^{n}(u - z_i)^{h_i}\)。将 \(\operatorname{Ker}f(u)\) 看作 \(\mathbf{C}[\mathbf{X}]\)-模，其外作用律为 \((h,x)\mapsto h(u)x\)，其中 \(h\in \mathbf{C}[\mathbf{X}]\)、\(x\in \operatorname{Ker}f(u)\)，可见 \(\operatorname{Ker} f(u)\) 是各 \(\operatorname{Ker}(u - z_i)^{h_i}\) 的直和，应用 Algebra, Chapter VII, § 2, no. 1, Proposition 1 即得。

\subsection{在线性表示中的应用}

命题 13. 设 G 是连通李群，\(\rho\) 是 G 在完备可赋范空间 E 上的解析线性表示。令 \(\mathrm{E}_1, \mathrm{E}_2\) 为 E 的两个闭向量子空间，满足 \(\mathrm{E}_2 \subset \mathrm{E}_1\)。以下条件等价：

\begin{enumerate}
\def\labelenumi{(\roman{enumi})}
\item
  \(\rho(g)x \equiv x \pmod{\mathbf{E}_2}\)，对所有 \(g \in \mathbf{G}\) 和所有 \(x \in \mathbf{E}_1\)；
\item
  \(\mathrm{L}(\rho)(\mathrm{L}(\mathrm{G}))\) 把 \(E_{1}\) 映入 \(E_{2}\)。
\end{enumerate}

\[
\begin{array}{l l l} \rho (g) x \equiv x & (\mathrm{mod} \mathrm{E} _ {2}) & \text {for all \(g\in G\) and all \(x\in E_{1}\)} \\ \Leftrightarrow \rho (\exp a) x \equiv x & (\mathrm{mod} \mathrm{E} _ {2}) & \text {for all \(a\in L(G)\) and all \(x\in E_{1}\)} \\ \Leftrightarrow (\exp \mathrm{L} (\rho) a) x \equiv x & (\mathrm{mod} \mathrm{E} _ {2}) & \text {for all \(a\in L(G)\) and all \(x\in E_{1}\)}. \end{array}
\]

另一方面，若 \(u \in \mathcal{L}(\mathrm{E})\)，则

\[
\begin{array}{r l} & \exp (\lambda u) x \equiv x \pmod {\mathrm{E} _ {2}} \quad \text {(for all} \lambda \in \mathrm{K} \text {and all} x \in \mathrm{E} _ {1} \\ & \Leftrightarrow u (\mathrm{E} _ {1}) \subset \mathrm{E} _ {2} \end{array}
\]

从而命题得证。

推论 1. \(\mathbf{E}_1\) 在 \(\rho\) 下稳定的充要条件是 \(\mathbf{E}_1\) 在 \(\mathrm{L}(\rho)\) 下稳定。

在命题 13 中取 \(E_{1} = E_{2}\) 即可。

推论 2. 假设 \(\rho\) 是有限维的。\(\rho\) 为单的（分别为半单的）的充要条件是 \(\mathbf{L}(\rho)\) 为单的（分别为半单的）。

这由推论 1 得出。

推论 3. 设 \(x \in E\)。\(x\) 在 \(\rho(G)\) 下不变的充要条件是 \(x\) 被 \(L(\rho)(L(G))\) 消去（即，按 Chapter I, § 3, Definition 3 的意义，\(x\) 在 \(L(\rho)\) 下不变）。

在命题 13 中取 \(\mathbf{E}_1 = \mathbf{K}x\) 和 \(\mathbf{E}_2 = 0\) 即可。

推论 4. 设 \(\rho'\) 是 \(G\) 在完备可赋范空间 \(E'\) 上的另一个解析线性表示。令 \(T \in \mathcal{L}(E, E')\)。以下条件等价：

\begin{enumerate}
\def\labelenumi{(\roman{enumi})}
\item
  \(\mathrm{T}\rho (g) = \rho '(g)\mathrm{T}\)，对所有 \(g\in \mathbf{G}\)
\item
  \(\mathrm{TL}(\rho)(a) = \mathrm{L}(\rho')(a)\mathrm{T}\)，对所有 \(a\in \mathrm{L}(\mathrm{G})\)
\end{enumerate}

令 \(\sigma\) 为 \(G\) 在 \(\mathcal{L}(E, E')\) 上的线性表示，它由 \(\rho\) 和 \(\rho'\) 导出（\(\S 3\)，no. 11, Corollary 1 to Proposition 41）。条件 (i) 意味着 \(T\) 在 \(\sigma(G)\) 下不变。条件 (ii) 意味着 \(T\) 被 \(L(\sigma)(L(G))\) 消去。因此只需应用推论 3。

推论 5. 假设 \(\rho\) 和 \(\rho'\) 是有限维的。\(\rho\) 和 \(\rho'\) 等价的充要条件是 \(\mathbf{L}(\rho)\) 和 \(\mathbf{L}(\rho')\) 等价。

这是推论 4 的特殊情形。

推论 6. 假设 \(\mathbf{G}\) 是有限维的。令 \(t \in \mathrm{U}(\mathrm{G})\)。\(\mathbf{L}_t\) 为右不变（分别地，\(\mathbf{R}_t\) 为左不变）的充要条件是 \(t\) 属于 \(\mathrm{U}(\mathrm{G})\) 的中心。

\(\mathrm{L}_t\) 为右不变（分别地，\(\mathbf{R}_t\) 为左不变）的充要条件是 \(\varepsilon_g * t = t * \varepsilon_g\)，对所有 \(g \in G\)，也就是 \((\operatorname{Int} g)_* t = t\)。存在一个整数 \(n\)，使得 \(t \in \mathrm{U}_n(G)\)。由 § 3, no. 12 的 Corollary 3 和 Proposition 45，\((\operatorname{Int} g)_* t = t\) 对所有 \(g \in G\) 成立，当且仅当 \([a, t] = 0\) 对所有 \(a \in \mathrm{L}(G)\) 成立，也就是当且仅当 \(t\) 与 \(\mathrm{U}(G)\) 交换。

\subsection{正规积分子群}

引理 3. 设 \(\mathbf{G}\) 是李群，\(\mathrm{H}_1\) 和 \(\mathrm{H}_2\) 是拓扑有可数基的积分子群，且 \(g \in \mathbf{G}\)。则

\[
g \mathrm{H} _ {1} g ^ {- 1} = \mathrm{H} _ {2} \Leftrightarrow (\operatorname{Ad} g) (\mathrm{L} (\mathrm{H} _ {1})) = \mathrm{L} (\mathrm{H} _ {2}).
\]

\(\operatorname{Ad} g = \mathrm{T}_{e}(\operatorname{Int} g)\)。因此，由结构传递，\((\operatorname{Int} g)(\mathrm{H}_{1})\) 的李代数为 \((\operatorname{Ad} g)(\mathrm{L}(\mathrm{H}_{1}))\)。由于 \(\mathrm{H}_{1}\) 和 \(\mathrm{H}_{2}\) 都有可数基，说集合 \(\mathrm{H}_{2}\) 与 \((\operatorname{Int} g)(\mathrm{H}_{1})\) 相等，等价于说积分子群 \(\mathrm{H}_{2}\) 与 \((\operatorname{Int} g)(\mathrm{H}_{1})\) 相等（no. 2, Corollary 3 to Proposition 3）。于是引理由定理 2 (i) 得出。

命题 14. 设 G 是李群，H 是拓扑有可数基的积分子群。以下条件等价：

\begin{enumerate}
\def\labelenumi{(\roman{enumi})}
\item
  H 是 G 的正规子群；
\item
  L(H) 在 Ad(G) 下不变。
\end{enumerate}

若进一步假设 G 连通，则这些条件还等价于：

\begin{enumerate}
\def\labelenumi{(\roman{enumi})}
\setcounter{enumi}{2}
\tightlist
\item
  L(H) 是 L(G) 的理想。
\end{enumerate}

若进一步假设 \(\mathbf{G}\) 单连通且 \(\mathbf{L}(\mathbf{H})\) 在 \(\mathbf{L}(\mathbf{G})\) 中具有有限余维，则这些条件蕴含 \(\mathbf{H}\) 是 \(\mathbf{G}\) 的李子群，且 \(\mathbf{G} / \mathbf{H}\) 单连通。

(i) \(\Leftrightarrow\) (ii) 的等价性由引理 3 得出。若进一步假设 G 连通，则条件 (ii) 等价于 L(H) 在 ad L(G) 下稳定（no. 5, Corollary 1 to Proposition 13 and § 3, no. 12, Proposition 44）。

假设 G 单连通且 L(H) 是 L(G) 中有限余维的理想。由第 3 号的 Theorem 3，存在单连通李群 G'，使 L(G') 同构于 L(G)/L(H)。存在一个从 L(G) 到 L(G') 的连续满态射 f，核为 L(H)。由第 1 号的 Theorem 1，存在从 G 到 G' 的态射 \(\phi\)，使 L( \(\phi\) ) = f.~该态射是下沉映射，因而其核 N 是 G 的李子群，且 L(N) = Ker f = L(H)。所以 H 是 N 的单位元分支，因而是 G 的李子群。令 \(\psi\) 为将 \(\phi\) 通过商得到的从 G/H 到 G' 的态射。该态射是 étale 的；由于 G 单连通，\(\psi\) 是从 G/H 到 G' 的同构。

推论 1. 设 \(G\) 是有限维单连通李群。令 \(m, h\) 为 \(L(G)\) 的李子代数，使得 \(L(G)\) 是 \(m\) 与 \(h\) 的半直积。令 \(M, H\) 为 \(G\) 中相应的积分子群。则 \(M\) 和 \(H\) 是 \(G\) 的单连通李子群，并且作为李群，\(G\) 是 \(M\) 与 \(H\) 的半直积。

由命题 14，H 是 G 的正规李子群，且李群 G/H 单连通。令 \(\pi\) 为从 G 到 G/H 的规范态射。存在从 G/H 到 M 的态射 \(\theta\)，使 L(θ) 是从 L(G)/L(H) = L(G)/h 到 L(M) = m 的规范同构。于是

\[
\mathrm{L} (\pi \circ \theta) = \mathrm{L} (\pi) \circ \mathrm{L} (\theta) = \mathrm{Id} _ {\mathrm{L} (\mathrm{G} / \mathrm{H})}
\]

从而 \(\pi \circ \theta = \mathrm{Id}_{\mathrm{G}/\mathrm{H}}\)。由第 1 号 Proposition 1 的 Corollary 1，\(\theta (\mathrm{G / H}) = \mathbf{M}\)。由 § 1, no. 4 的 Proposition 8，M 是 G 的李子群，且李群 G 是 M 与 H 的半直积。

推论 2. 设 G 是有限维单连通李群，H 是 G 的正规连通李子群，\(\pi\) 是从 G 到 G/H 的规范态射。

\begin{enumerate}
\def\labelenumi{(\roman{enumi})}
\item
  存在从 G/H 到 G 的解析映射 \(\rho\)，使得 \(\pi \circ \rho = \mathrm{Id}_{\mathrm{G / H}}\)。
\item
  对每个具有 (i) 中性质的映射 \(\rho\)，映射 \((h, m) \mapsto h\rho(m)\) 从 \(\mathbf{H} \times (\mathbf{G}/\mathbf{H})\) 到 \(\mathbf{G}\)，是解析流形同构。
\item
  H 和 G/H 都是单连通的。
\end{enumerate}

令 \(n = \dim G - \dim H\)。当 \(n = 0\) 时推论显然成立。我们对 \(n\) 作归纳论证。

假设存在一个包含 L(H) 且不同于 L(G) 和 L(H) 的 L(G) 理想。令 H' 为相应的 G 的连通李子群。令 \(\pi_{1}: G \to G/H'\) 和 \(\pi_{2}: H' \to H'/H\) 为规范态射。由归纳假设，存在解析映射 \(\rho_{1}: G/H' \to G\)、\(\rho_{2}: H'/H \to H'\)，使得 \(\pi_{1} \circ \rho_{1} = Id_{G/H'}\)、\(\pi_{2} \circ \rho_{2} = Id_{H'/H}\)。令

\[
\pi_ {3} \colon \mathrm{G} / \mathrm{H} \rightarrow \mathrm{G} / \mathrm{H} ^ {\prime}
\]

为规范态射。若 \(x \in \mathrm{G} / \mathrm{H}\)，且 \(y\) 是 \(x\) 在 \(\mathrm{G}, y\) 中的代表元，则 y 和 \(\rho_1(\pi_3(x))\) 模 \(\mathrm{H}'\) 有相同的规范像，因而 \(x^{-1}\pi (\rho_1(\pi_3(x))) \in \mathrm{H}' / \mathrm{H}\)。令

\[
\rho (x) = \rho_ {1} (\pi_ {3} (x)) \rho_ {2} (\pi (\rho_ {1} (\pi_ {3} (x))) ^ {- 1} x) \in G.
\]

显然，\(\rho\) 是从 G/H 到 G 的解析映射，并且

\[
\begin{array}{r l} \pi (\rho (x)) & = \pi (\rho_ {1} (\pi_ {3} (x))) \pi_ {2} (\rho_ {2} (\pi (\rho_ {1} (\pi_ {3} (x))) ^ {- 1} x)) \\ & = \pi (\rho_ {1} (\pi_ {3} (x))) \pi (\rho_ {1} (\pi_ {3} (x))) ^ {- 1} x = x. \end{array}
\]

如果现在 \(\mathbf{L}(\mathbf{G})\) 的理想中，包含 \(\mathbf{L}(\mathbf{H})\) 的只有 \(\mathbf{L}(\mathbf{G})\) 和 \(\mathbf{L}(\mathbf{H})\)，则李代数 \(\mathrm{L(G) / L(H)}\) 要么是一维的，要么是单的。在这两种情形下，\(\mathrm{L(G)}\) 都是某个子代数与 \(\mathrm{L(H)}\) 的半直积；第一种情形显然，第二种情形由 Chapter I, § 6, Corollary 3 to Theorem 5 得出。因此断言 (i) 由推论 1 得出。

断言 (ii) 显然。断言 (iii) 由 (i) 和 (ii) 得出。

当 G 是无限维的或 H 不是正规子群时，推论 2 的结论不再必然成立（练习 8 和 15）。

推论 3. 设 \(\mathbf{G}\) 是有限维连通李群，\(\mathbf{H}\) 是 \(\mathbf{G}\) 的正规连通李子群。从 \(\pi_1(\mathbf{H})\) 到 \(\pi_1(\mathbf{G})\) 的规范态射是单射。

令 \(G_{1}\) 为 G 的万有覆盖，\(\lambda\) 为从 \(G_{1}\) 到 G 的规范映射。将 \(G_{1}\) 的李代数与 L(G) 识别。\(\lambda^{-1}(H)\) 是 \(G_{1}\) 的李子群，在 \(G_{1}\) 中正规，且其李代数为 L(H)。令 \(H_{1}\) 为 \(\lambda^{-1}(H)\) 的单位元分支，\(\lambda_{1} = \lambda | H_{1}\)。则 \(\mathrm{L}(H_{1}) = \mathrm{L}(H)\)，从而 \(\lambda_{1}\) 是从 \(H_{1}\) 到 H 的 étale 态射。另一方面，\(H_{1}\) 是单连通的（推论 2），因而可与 H 的万有覆盖识别。于是根据 General Topology, Chapter XI，从 \(\pi_{1}(H)\) 到 \(\pi_{1}(G)\) 的规范态射可识别为从 Ker \(\lambda_{1}\) 到 Ker \(\lambda\) 的规范嵌入。

\subsection{取值于李代数的微分形式的原函数}

命题 15. 设 \(\mathrm{G}\) 是李群，\(\mathrm{M}\) 是 \(\mathrm{C}^r (r\geqslant 2)\) 类流形，\(\alpha\) 是 \(\mathrm{C}^{r - 1}\) 类的 1 次微分形式，定义在 \(\mathrm{M}\) 上且取值于 \(\mathrm{L}(\mathrm{G})\)，满足 \(d\alpha +[\alpha ]^2 = 0\)。假设 \(\mathrm{M}\) 单连通。对所有 \(x\in \mathrm{M}\) 和 \(s\in \mathrm{G}\)，存在唯一的映射 \(f\)，其属于 \(\mathrm{C}^{r - 1}\) 类，定义在 \(\mathrm{M}\) 上且取值于 \(\mathrm{G}\)，使得 \(f(x) = s\) 且 \(f^{-1}.df = \alpha\)。

\(f\) 的唯一性由 § 3, no. 17, Corollary 2 to Proposition 59 以及 \(\mathbf{M}\) 连通这一事实得出。下面证明 \(f\) 的存在性。存在一个开覆盖 \((\mathrm{U}_i)_{i\in \mathbf{I}}\)，它覆盖 \(\mathbf{M}\)，并且对所有 \(i\in \mathbf{I}\)，存在映射 \(g_i: \mathrm{U}_i\to \mathrm{G}\)，它属于 \(\mathrm{C}^{r - 1}\) 类，并且 \(g_i^{-1}.d g_i = \alpha\) 在 \(\mathrm{U}_i\) 上成立（§ 4, no. 6, Theorem 5）。由 § 3, no. 17, Corollary 2 to Proposition 59，\(g_i g_j^{-1}\) 在 \(\mathrm{U}_i\cap \mathrm{U}_j\) 上局部为常值。令 \(g_i g_j^{-1} = g_{ij}\)。令 \(\mathbf{G}_d\) 为带离散拓扑的群 \(\mathbf{G}\)。映射 \(g_{ij}: \mathrm{U}_i\cap \mathrm{U}_j\to \mathrm{G}_d\) 连续，并且 \(g_{ij}g_{jk} = g_{ik}\) 在 \(\mathrm{U}_i\cap \mathrm{U}_j\cap \mathrm{U}_k\) 上成立。由于 \(\mathbf{M}\) 单连通，存在连续映射 \(\lambda_i: \mathrm{U}_i\to \mathrm{G}_d\)，使得 \(g_i g_j^{-1} = \lambda_i\lambda_j^{-1}\) 在 \(\mathrm{U}_i\cap \mathrm{U}_j\) 上成立。令 \(g\) 为从 \(\mathbf{M}\) 到 \(\mathbf{G}\) 的映射，其在 \(\mathrm{U}_i\) 上的限制为 \(\lambda_i^{-1}g_i\)，对所有 \(i\in \mathbf{I}\) 均如此。这个映射属于 \(\mathrm{C}^{r - 1}\) 类，且 \(g^{-1}d g = \alpha\)。映射 \(f\) 从 \(\mathbf{M}\) 到 \(\mathbf{G}\)，由 \(f = s(g(x))^{-1}g\) 定义，并满足条件 \(f^{-1}.df = \alpha\) 和 \(f(x) = s\)。

\subsection{从无穷小作用律到作用律的过渡}

引理 4. 设 \(\mathbf{G}\) 是连通拓扑群，\(\mathbf{X}\) 是 Hausdorff 拓扑空间，\(f_{1}, f_{2}\) 是 \(\mathbf{G}\) 在 \(\mathbf{X}\) 上的左（分别为右）作用律，使得对所有 \(x \in \mathbf{X}\)，从 G 到 X 的映射 \(s \mapsto f_1(s, x)\)、\(s \mapsto f_2(s, x)\) 连续。假设存在 \(\{e\} \times X\) 在 G \(\times\) X 中的某个邻域 V，使 \(f_1\) 和 \(f_2\) 在 V 上重合。则 \(f_1 = f_2\)。

令 \(x \in X\)，A 为由 \(g \in G\) 且满足 \(f_{1}(g, x) = f_{2}(g, x)\) 的元素组成的集合。则 A 在 G 中是闭集。另一方面，令 \(g \in A\)；记 \(y = f_{1}(g, x) = f_{2}(g, x)\)。存在 G 中 e 的一个邻域 U，使得 \(f_{1}(t, y) = f_{2}(t, y)\) 对 \(t \in U\) 成立，换言之，使得 \(f_{1}(t', x) = f_{2}(t', x)\) 对 \(t' \in Ug\)（分别为 gU）成立。因此 A 在 G 中是开集，故 A = G。

命题 16. 设 G 是连通李群，X 是 \(\mathbf{C}^r\) 类 Hausdorff 流形，\(f_1, f_2\) 是 G 在 X 上的 \(\mathbf{C}^r\) 类左（分别为右）作用律。若与 \(f_1\) 和 \(f_2\) 相关的无穷小作用律相等，则 \(f_1 = f_2\)。

由 § 4, no. 7, Corollary to Proposition 11，存在 \(\{e\} \times \mathbf{X}\) 在 \(\mathbf{G} \times \mathbf{X}\) 中的一个邻域 V，使得 \(f_{1}\) 和 \(f_{2}\) 在 V 上重合。因此 \(f_{1} = f_{2}\)（引理 4）。

引理 5. 设 G 是单连通拓扑群，X 是 Hausdorff 拓扑空间，U 是 G 中 e 的一个开邻域，\(\psi\) 是从 U × X 到 X 的连续映射，满足 \(\psi(e, x) = x\)，并且 \(\psi(s, \psi(t, x)) = \psi(st, x)\)，对所有 \(x \in X\) 以及 U 中满足 \(st \in U\) 的 s, t。令 W 为 e 的一个对称连通开邻域，使 \(W^{3} \subset U\)。存在唯一一个 G 在 X 上的连续左作用律 \(\psi'\)，使 \(\psi'\) 和 \(\psi\) 在 W × X 上重合。若 G 是李群，X 是 \(\mathbf{C}^r\) 类流形且 \(\psi\) 属于 \(\mathbf{C}^r\) 类，则 \(\psi'\) 属于 \(\mathbf{C}^r\) 类。

由引理 4，\(\psi\) 的唯一性成立。令 P 为 X 的置换群。对于 \(u \in W^{3}\)，映射 \(x \mapsto \psi(u, x)\) 是 P 的元素 \(f(u)\)，并且

\[
f \left(u _ {1} u _ {2} u _ {3}\right) = f \left(u _ {1}\right) f \left(u _ {2}\right) f \left(u _ {3}\right)
\]

对 W 中的 \(u_{1}, u_{2}, u_{3}\)。应用第 1 号引理 1，得到从 G 到 P 的态射 \(f'\)，它延拓 \(f|\mathrm{W}\)。令 \(\psi'(g,x)=f'(g)(x)\)，对所有 \((g,x)\in\mathbf{G}\times\mathbf{X}\)。则 \(\psi'\) 是 G 在 X 上的左作用律，并在 W × X 上与 \(\psi\) 重合。由于 \(\psi'(g,\psi'(g',x))=\psi'(gg',x)\) 对 \((g,g',x)\in\mathbf{G}\times\mathbf{G}\times\mathbf{X}\) 成立，\(\psi\) 在 W × X 上的连续性蕴含 \(\psi'\) 在对所有 \(g\in\mathbf{G}\) 的 gW × X 上连续。因此 \(\psi'\) 连续。若 \(\psi\) 属于 \(\mathbf{C}^r\) 类，同样可见 \(\psi'\) 属于 \(\mathbf{C}^r\) 类。

定理 5. 设 G 是单连通李群，X 是 \(\mathbf{C}^r\) 类紧流形（\(r \geqslant 2\)），\(a \mapsto \mathrm{D}_a\) 是 L(G) 在 X 上的 \(\mathbf{C}^{r-1}\) 类左（分别为右）无穷小作用律。存在唯一一个 G 在 X 上的 \(\mathbf{C}^{r-1}\) 类左（分别为右）作用律，使其相应的无穷小作用律为 \(a \mapsto \mathrm{D}_a\)。

唯一性由命题 16 得出。由 § 4, no. 7, Corollary 1 to Theorem 6，存在 \(\{e\} \times \mathbf{X}\) 在 \(\mathrm{G} \times \mathrm{X}\) 中的一个邻域 V，以及 G 在 X 上定义于 V 的 \(\mathrm{C}^{r-1}\) 类左（分别为右）作用律片段，使其相应的无穷小作用律为 \(a \mapsto \mathrm{D}_a\)。由于 X 紧，可以设 V 形如 \(U \times X\)，其中 U 是 G 中 e 的一个开邻域。于是只需应用引理 5。
